\newif\ifdark
  \darkfalse
  %\darktrue % triggers dark coloured pdf for late night work

\newif\iffutureversion
  \futureversiontrue
  %\futureversionfalse

\newif\ifdraftversion
  %\draftversionfalse
  \draftversiontrue

\newif\iffullversion
  % \fullversionfalse % use for camera-ready and author version only
  \fullversiontrue 

\newif\ifauthversion
  \authversiontrue % triggers copyright footnote for non-full eprint author version
  % \authversionfalse % triggers submission version for full version (anonymous, suppl material)

\newif\iftutorial % to flag parts that should be excluded in a SoK
  \tutorialtrue
  %\tutorialfalse

%%% BEGINNING OF DO NOT EDIT ZONE %%%

\newif\ifsubmission % best not to manually change this flag
  \iffullversion
    \ifauthversion
      \submissionfalse
    \else
      \submissiontrue
    \fi
  \else
    \submissionfalse
\fi

\newif\ifeprintfull % best not to manually change this flag
  \iffullversion
    \ifauthversion
      \eprintfulltrue
    \else
      \eprintfullfalse
    \fi
  \else
    \eprintfullfalse
  \fi

\iffullversion \ifdraftversion \documentclass[runningheads,oneside]{llncs}
               \else           \documentclass[runningheads,oneside      ]{llncs} \fi
\else          \ifdraftversion \documentclass[runningheads,            ]{llncs}
               \else           \documentclass[runningheads,                  ]{llncs} \fi\fi

%%% END OF DO NOT EDIT ZONE %%%
                
% HH: Added to make table-of-contents work
%% see: https://tex.stackexchange.com/questions/272200/table-of-contents-showing-the-title-as-only-entry-latex
  \setcounter{tocdepth}{1} % Exclude subsections
  \makeatletter
  \renewcommand*\l@author[2]{}
  \renewcommand*\l@title[2]{}
  \makeatother
%%

\usepackage{cite}
\usepackage{bbm} % for thick 1. wait, isn't there already a package in cryptocodeh for this?
\usepackage{dashbox} 

%% For mono-width matrices
  \usepackage{nicematrix}
  %% grid matrix; scales all column widths to fit the widest 
  \NewDocumentEnvironment{gmatrix}{O{auto}} % default = auto, or pass a fixed width
    {\begin{bNiceMatrix}[columns-width=#1]}
    {\end{bNiceMatrix}}
  %% For square matrices to actually render as square
  \newlength{\sqw}
  \NewDocumentEnvironment{sqmatrix}{O{1}}{%
    \setlength{\sqw}{#1\baselineskip}%   cell width = (multiple of) row height
    \setlength{\arraycolsep}{0pt}%        kill the inter-column gap
    \begin{bNiceMatrix}[columns-width=\sqw]%
  }{\end{bNiceMatrix}}

\ifeprintfull
  \ifdraftversion
    %\usepackage[reset,a4paper]{geometry}
    %\geometry{textwidth=12.2cm, textheight=19.3cm}
    \usepackage[a4paper,margin=1.0in]{geometry}
  \else
    \usepackage[a4paper,margin=1.0in]{geometry}
  \fi

  \usepackage[section]{placeins} % Gets overwritten by ! tag in float
\fi

\ifsubmission
  \usepackage[reset]{geometry}
  \geometry{textwidth=12.2cm, textheight=19.3cm}
  \pagestyle{plain}
\else
  \pagestyle{headings}
\fi

\usepackage[svgnames,table]{xcolor}

%% Hacks to enable grayed-out ToC entries %%
  \definecolor{draftgrey}{gray}{0.55}
  \newcommand{\plannedsection}[1]{%
    \stepcounter{section}%
    \addtocontents{toc}{%
      \protect\begingroup\protect\color{draftgrey}%
      \protect\contentsline{section}{\protect\numberline{\thesection}#1}{\textcolor{draftgrey}{
        %/ % symbol used in place of a page number
      }}{}%
      \protect\endgroup}%
  }
  % dispatcher: real content when building the future version, phantom otherwise
  \newcommand{\futuresect}[2]{% {file}{ToC title}
    \iffutureversion \input{#1}\clearpage \else \plannedsection{#2}\fi
  }

% Todonotes
\ifdraftversion
  \usepackage{changes} % also loads todonotes
\else
  \usepackage[final,todonotes=disable]{changes}
\fi

% Things Kat has added
  \usepackage{nicodemus}  % For easier handling of pseudocode

% Things Hans has added
  \usepackage{slashed}
  \usepackage[normalem]{ulem}
  \usepackage{booktabs}
  \usepackage{nicefrac}
  \usepackage{braket}

% thm-restate: lets a lemma/theorem be stated once (in a \begin{restatable}
% block) and replayed verbatim with the SAME number elsewhere (e.g. an
% appendix) via a generated \<macro>* command. Works directly on llncs's
% \spnewtheorem environments; no \declaretheorem redefinition needed.
% Loaded before cryptocode so amsthm (pulled in by thmtools) precedes it;
% hyperref/cleveref still load later, as cleveref requires.
\usepackage{thmtools,thm-restate}

\input{./cryptocodeh/cryptocodeh}
\input{./cryptocodeh/cryptocrakz}
\graphicspath{{../figures/}}

\usepackage{orcidlink}
\renewcommand{\orcidID}[1]{\orcidlink{#1}}

\usepackage{array}
\newcolumntype{C}{>{$}c<{$}}

% Dark mode: Colour scheme like a moonlit night or something.
\ifdark
	 \pagecolor[rgb]{0.05,0.16,0.21}
	 \color[rgb]{0.53,0.58,0.59}
\fi

\newcommand{\crakzfig}[2]{
  \begin{figure}[t!]
    \centering
    \includegraphics[keepaspectratio=true]{./crakz/#1}
    \caption{#2}
    {\label{fig:#1}}
  \end{figure}
}

\newcommand{\crakzfigcode}[2]{
  \begin{figure}[t!]
    \centering
    \input{./crakz/#1}
    \caption{#2}
    {\label{fig:#1}}
  \end{figure}
}

%%% Commands go here

%%% Arrows
\newcommand{\apply}[1]{\overset{#1}{\longrightarrow}}

%%% Sets and spaces
\newcommand{\funcspace}[2]{\pcsetstyle{F}_{#1 \to #2}}

%%% Game interface
\renewcommand{\pcthen}{:}

%%% Oracles
\newcommand{\ro}{\hash}

%%% Algorithms
\newcommand{\randfunc}{\ensuremath{\mathcal{R}}}
\newcommand{\unitary}{\ensuremath{\pcalgostyle{U}}}

%%% Quantum objects
\newcommand{\qubit}{\mathsf{q}}
\newcommand{\register}[1][R]{\ensuremath{\mathbb{#1}}}
\newcommand{\qcomp}{\register[Q]}

%%% Classical objects
\newcommand{\prfkey}{\ensuremath{k}}

%%% Gates
\renewcommand{\unitary}[1][U]{\mathsf{#1}}
\newcommand{\Xgate}{\unitary{X}}
\newcommand{\Zgate}{\unitary{Z}}
\newcommand{\Hgate}{\unitary{H}}
\newcommand{\CXgate}{\unitary{CX}}

\newcommand{\PKE}{\ensuremath{\mathsf{PKE}}\xspace}
\newcommand{\DPKE}{\ensuremath{\mathsf{DPKE}}\xspace}
\newcommand{\PG}{\ensuremath{\mathsf{PmGen}}\xspace}
\newcommand{\KG}{\ensuremath{\mathsf{KeyGen}}\xspace}
\newcommand{\Enc}{\ensuremath{\mathsf{Enc}}\xspace}
\newcommand{\Dec}{\ensuremath{\mathsf{Dec}}\xspace}

\newcommand{\KEM}{\ensuremath{\mathsf{KEM}}\xspace}
\newcommand{\Encaps}{\ensuremath{\mathsf{Encaps}}\xspace}
\newcommand{\Decaps}{\ensuremath{\mathsf{Decaps}}\xspace}
\newcommand{\KEMoutputKey}{\ensuremath{\mathsf{k}}\xspace}

% PKE /KEM spaces
\newcommand{\PKSpace}{\ensuremath{\mathcal{PK}}\xspace}
\newcommand{\SKSpace}{\ensuremath{\mathcal{SK}}\xspace}
\newcommand{\MSpace}{\ensuremath{\mathcal{M}}\xspace}
\newcommand{\KSpace}{\ensuremath{\mathcal{K}}\xspace}
\newcommand{\CSpace}{\ensuremath{\mathcal{C}}\xspace}
\newcommand{\RSpace}{\ensuremath{\mathcal{R}}\xspace}

% Security notions for PKEs/KEMs
\mathchardef\mhyphen="2D % Define a "math hyphen"
\newcommand{\OW}{\ensuremath{\mathsf{OW}}\xspace}
\newcommand{\IND}{\ensuremath{\mathsf{IND}}\xspace}
\newcommand{\CPA}{\ensuremath{\mathsf{CPA}}\xspace}
\newcommand{\CCA}{\ensuremath{\mathsf{CCA}}\xspace}
\newcommand{\OWCPA}{\ensuremath{\OW\mhyphen\CPA}\xspace}
\newcommand{\INDCPA}{{\ensuremath{\IND\mhyphen\CPA}}\xspace}
\newcommand{\INDCCA}{{\ensuremath{\IND\mhyphen\CCA}}\xspace}

% === standard FO notation
\newcommand{\T}{\ensuremath{\mathsf{T}}\xspace}
\newcommand{\TEnc}{\ensuremath{\Enc'}\xspace}
\newcommand{\TDec}{\ensuremath{\Dec'}\xspace}

\newcommand{\U}{\ensuremath{\mathsf{U}}\xspace}
\newcommand{\notp}{{\not\perp}} % <- for implicit rejection
%\newcommand{\UExp}{\ensuremath{\U^\perp}\xspace}
%\newcommand{\UExpMess}{\ensuremath{\UExp_{m}}\xspace}
\newcommand{\UImp}{\ensuremath{\U^\notp}\xspace}
\newcommand{\UImpMess}{\ensuremath{\UImp_{m}}\xspace}

%\newcommand{\KEMExp}{\ensuremath{\KEM^{\perp}}\xspace}
%\newcommand{\KEMExpMess}{\ensuremath{\KEMExp_m}\xspace}
\newcommand{\KEMImp}{\ensuremath{\KEM^{\notp}}\xspace}
\newcommand{\KEMImpMess}{\ensuremath{\KEMImp_m}\xspace}

%\newcommand{\KGExp}{\ensuremath{\KG}\xspace}
\newcommand{\KGImp}{\ensuremath{\KG^\notp}\xspace}
\newcommand{\rejectionSeed}{\ensuremath{z}\xspace}
\newcommand{\EncapsMess}{\ensuremath{\Encaps_m}\xspace}
%\newcommand{\DecapsExpMess}{\ensuremath{\Decaps^\perp_m}\xspace}
\newcommand{\DecapsImpMess}{\ensuremath{\Decaps^\notp_m}\xspace}

\newcommand{\FO}{\ensuremath{\mathsf{FO}}\xspace}
%\newcommand{\FOExp}{\ensuremath{\FO^{\perp}}\xspace}
%\newcommand{\FOImp}{\ensuremath{\FO^\notp}\xspace}
%\newcommand{\FOMess}{\ensuremath{\FO_m}\xspace}
%\newcommand{\FOMessExp}{\ensuremath{\FOExp_m}\xspace}
%\newcommand{\FOMessImp}{\ensuremath{\FOImp_m}\xspace}


% Hash functions used by FO
\newcommand{\hashTS}[1]{\ensuremath{\mathsf{#1}}\xspace} % <- for typesetting hash functions
\newcommand{\ROkey}{\hashTS{H}}
\newcommand{\ROrand}{\hashTS{G}}
\newcommand{\ROimp}{\hashTS{H'}}


% Stuff for security definitions
\newcommand{\AdversaryA}{\ensuremath{\mathcal{A}}\xspace}
\newcommand{\Advantage}{\ensuremath{\mathrm{Adv}}\xspace}
\newcommand{\oracleDecaps}{\ensuremath{\mathsc{Dec}}\xspace}
\newcommand{\bool}[1]{\ensuremath{\llbracket #1\rrbracket}\xspace}
%\newcommand{\sample}{{\:\leftarrow\hspace*{-3pt}\raisebox{.5pt}{$\scriptscriptstyle\$$}\:}}

% most of the commands are still in nicodemus.sty. However, e.g. texstudio will not allow for autocompletion of commands in nicodemus.sty :-/
\newenvironment{nicodemus}[1][\thenicolinenr]{%  % parameter controls initial line number
	\begin{enumerate}[
		topsep=0ex,
		label=\nicolinenrformat\PaddingUp*,
		ref=\nicorefprefix\PaddingUp*,
		align=right,
		leftmargin=0em,
		itemindent=!,
		labelindent=0em,
		labelwidth=\nicolinenrwidth,
		labelsep=\nicolinenrsep,
		listparindent=\parindent,
		noitemsep,
		]%
		\setcounter{enumi}{#1}%
		\addtocounter{enumi}{-1}%
	}{%
	\end{enumerate}%
	\addtocounter{enumi}{1}%
	\setcounter{nicolinenr}{\theenumi}%
}
 % <- Kat's collection of macros for stuff on FO

% Commenting
%-------------------------------------------------------------------------------%

\definechangesauthor[name={Kathrin}, color=purple]{TS}
\definechangesauthor[name={Hans}, color=darkgreen]{HH}

\newcommand{\hhnote}[2][]{\todo[inline,color=Linen]{\textbf{Hans:} \if!#1! #2 \else (#1)\\ #2\fi}\xspace}
\newcommand{\khnote}[2][]{\todo[inline,color=violet,textcolor=white]{\textbf{Kathrin:} \if!#1! #2 \else (#1)\\ #2\fi}\xspace}
\newcommand{\revnote}[1]{\todo[inline]{\textbf{Reviewer:} #1}\xspace}

\newcommand\todoenum[2][]{\todo[inline,color=Cornsilk,caption={2do}]{
  \begin{minipage}{\textwidth-4pt}\textbf{#1}\begin{enumerate}#2\end{enumerate}\end{minipage}}}
\newcommand\doneenum[2][]{\todo[color=FloralWhite,caption={2do}]{
  \begin{minipage}{\textwidth-4pt}\textbf{#1}\begin{enumerate}#2\end{enumerate}\end{minipage}}}
\newcommand\oldtodoenum[2][]{\todo[inline,color=FloralWhite,caption={2do}]{
  \begin{minipage}{\textwidth-4pt}\textbf{#1}\begin{enumerate}#2\end{enumerate}\end{minipage}}}

%------------------------------------------------------------------------------%

\newcommand{\pagehack}{\ifsubmission\clearpage\fi}

%% Cross-referencing via cleveref (migrated from autoref).
% IMPORTANT: cleveref must load AFTER hyperref, which is pulled in by
% orcidlink above. llncs defines its theorem-like environments with
% \spnewtheorem (not \newtheorem), each with its own counter, so cleveref
% cannot auto-name them: every environment needs an explicit \crefname.
% Names are kept capitalized (matching the old *autorefname abbreviations),
% with \Crefname identical, so \cref and \Cref render the same.
\usepackage[nameinlink]{cleveref}
\crefname{section}{Sect.}{Sect.}            \Crefname{section}{Sect.}{Sect.}
\crefname{subsection}{Sect.}{Sect.}         \Crefname{subsection}{Sect.}{Sect.}
\crefname{subsubsection}{Sect.}{Sect.}      \Crefname{subsubsection}{Sect.}{Sect.}
\crefname{part}{Part}{Parts}                \Crefname{part}{Part}{Parts}
\crefname{figure}{Fig.}{Figs.}              \Crefname{figure}{Fig.}{Figs.}
\crefname{table}{Table}{Tables}             \Crefname{table}{Table}{Tables}
\crefname{theorem}{Thm.}{Thms.}             \Crefname{theorem}{Thm.}{Thms.}
\crefname{lemma}{Lemma}{Lemmas}             \Crefname{lemma}{Lemma}{Lemmas}
\crefname{definition}{Def.}{Defs.}          \Crefname{definition}{Def.}{Defs.}
\crefname{corollary}{Cor.}{Cors.}           \Crefname{corollary}{Cor.}{Cors.}
\crefname{example}{Example}{Examples}       \Crefname{example}{Example}{Examples}
% NB: the custom boxed exercise env (below) steps "exercisecounter".
\crefname{exercisecounter}{Exercise}{Exercises} \Crefname{exercisecounter}{Exercise}{Exercises}

%%% Boxed exercise environment + sub-exercises (a), b), ...).
% tcolorbox is loaded by cryptocodeh (input above); environ/enumitem here.
\usepackage{environ}
\usepackage{enumitem}
\newcounter{exercisecounter}
\renewenvironment{exercise}[1][]{%
  \refstepcounter{exercisecounter}%
  \begin{tcolorbox}[colback=white, colframe=black!50, boxrule=0.4pt, arc=0pt]
    \textbf{Exercise \theexercisecounter\ifx\\#1\\\else\ (#1)\fi.}
}{%
  \end{tcolorbox}%
}
% Sub-exercises: \ref yields e.g. "3a" for individual parts.
\newlist{subex}{enumerate}{1}
\setlist[subex]{label=\textbf{\alph*)}, leftmargin=*, ref=\theexercisecounter\alph*}
% \cpageref wording: "p. 17" / "pp. 17--19" instead of "page"/"pages".
\crefname{page}{p.}{pp.}                     \Crefname{page}{P.}{Pp.}
% Appendix reference helper (hyperref-based; appendix-specific "App." name).
\newcommand{\aref}[1]{\hyperref[#1]{App.~\ref*{#1}}}
% Range helper, now via cleveref (e.g. "Sects. 3 to 5").
\newcommand{\multiref}[2]{\crefrange{#1}{#2}}

% cryptocode + cleveref fix (cryptocode docs "8.3 Hyperref"). Superset of
% \pcfixhyperref: also adds \cref formats for pseudocode line numbers
% ("line 5" / "lines 3 to 7"). Installs an \AtBeginDocument hook, so it
% MUST be called here in the preamble, not after \begin{document}.
\pcfixcleveref

% \newcommand{\defer}[1]{Appendix~\ref{#1}}   %full version
\newcommand{\defer}[1]{\ref{#1}}              %submission
% \newcommand{\defer}[1]{the full version}      %camera ready

% Start of document
%-------------------------------------------------------------------------------%
\begin{document}
  \let\cleardoublepage\clearpage

  \title{
  \iffullversion
    The Quantum Random Oracle Toolbox
    % \thanks{
    %   An extended abstract of this paper will appear at xxx;
    %   this is the full version.
    % }
  \else
    The Quantum Random Oracle Toolbox
    \ifauthversion
    \thanks{
      \copyright IACR 2024. This article is the final version submitted by the authors
      to the IACR and to Springer-Verlag on <date>. The version published
      by Springer-Verlag is available at <DOI>.
    }
    \fi
  \fi
}

\titlerunning{
    The QROT
}

\ifsubmission
  \author{
    \vspace*{-5mm}
  }
  \institute{
    \vspace*{-5mm}
  }
\else
   \author{
    Hans Heum \orcidlink{0000-0003-0527-2999} %\and
  }
  \authorrunning{H.~Heum} 

  \institute{ 
    NTNU - Norwegian University of Science and Technology, Trondheim, Norway.\\
    \email{hans.heum@ntnu.no}
  }
\fi

\maketitle

\begin{abstract}
  Proofs of security in the Quantum Random Oracle Model (QROM) face three main
    challenges, relating to the non-applicability of several conceptually simple yet
    powerful proof techniques from its classical counterpart (ROM):

  \begin{enumerate}
    \item Recording: Or, how to extract QRO queries.
    \item Reprogramming: Or, how to alter specific outputs of the QRO.
    \item Rewinding: Or, how to derandomize quantum oracle algorithms.
  \end{enumerate}

  \noindent
  In this gentle introduction to QROM proof techniques, we will show how to
    deal with the first two. We will
    do this by taking a simple encryption scheme due to Bellare and
    Rogaway~\cite{CCS:BelRog93} as our motivating example, and prove increasingly
    strong statements for it as we move through the literature on quantum
    random oracles and gather stronger and stronger tools.
    %
    We will only assume knowledge of basic linear algebra and standard
    game-hopping techniques (though prior exposure to
    classical random oracles doesn't hurt).

  \begin{keywords}
    %Concrete Security \textperiodcentered~
    Quantum Random Oracles \textperiodcentered~
    One-Way to Hiding \textperiodcentered~
    Compressed Oracle Technique
  \end{keywords}
\end{abstract}

\ifeprintfull
  \newpage
  \tableofcontents
  \newpage
\fi
  % maketitle, abstract, ToC

  % Main body

  \section{Preface}
  \khnote{test}
  \hhnote{1st draft done}

  \noindent
  So you want to prove post-quantum security in the random oracle model.

  Proofs of security in the ROM are often conceptually
    simple, and yet regularly achieve strong security claims with
    near-tightness for highly efficient constructions---which is as close as
    one gets to a holy trinity for provable security. In contrast, proofs in the
    Quantum Random Oracle Model are often hard to parse without specialized
    knowledge in quantum information theory, and infamously lossy.

  And yet, if we want our proofs of security to have merit in a post-quantum
    world, it seems that our choice is between the Standard Model and the
    QROM. This is because random oracles model hash functions, which are
    publicly available algorithms that anyone can run locally on their machine.
    Therefore, there is nothing stopping someone with a quantum computer to
    boot it up and run it on a superposition of inputs. If we want to
    accurately model quantum attackers, then, our model must allow for such
    strategies. 

  In the standard model, this is not a problem at all,
    since then the hash functions in question is provided to the attacker
    as-is. But if QROM proofs are lossy, then Standard Model constructions
    that achieve the same goals as their ROM counterparts are often highly
    inefficient in comparison (and, as every cryptographer knows, efficiency
    tends to trump security in the real world).

  Therefore, as we inch ever closer to the looming threat of
    cryptographically-relevant quantum computers, it is important that the
    barrier for proving security in the QROM is lowered as much as
    possible, so that the community may together bring our transition to a
    post-quantum world into effect, in a safe, efficient, and (provably) secure
    manner.

  \paragraph{ROM Techniques Don't Transfer \ldots}
    One of the most useful facts about security reductions in the random oracle model
      is that anyone can efficiently \emph{simulate} a random oracle, simply by
      drawing the random outputs one-at-a-time and storing them in an
      appropriate list. This technique, known as \emph{lazy sampling}, allows
      the reduction to not only implement the oracle efficiently, but also to
      see what input values the adversary is interested in, \emph{and}
      adaptively reprogram the oracle on inputs of its choice.
      Indeed, much of the strength of the random oracle model lies in exactly
      these capabilities.

    None of this transfers directly to the case of quantum-accessible oracles.
      To see why, consider an adversary that queries the oracle on a
      superposition of all possible inputs at once: the simulator would have to
      store an exponentially large list of input values in order process just
      this single query!

  \paragraph{\ldots But the Concepts Do.}
    And yet, the past decade has seen the development of techniques that
      recover most of these capabilities, albeit at a loss. Prime among these
      techniques are the various One-Way to Hiding (O2H) lemmas, which allow
      for careful reprogramming of the quantum random oracle (as well as
      limited extraction of the input values); Quantum Online-Extractability
      (QOE), which as the name suggests allows for on-the-fly extraction of oracle inputs;
      and the Compressed Random Oracle (CRO) technique, which allows one to
      efficiently simulate a quantum random oracle, in direct
      analogy to classical ``lazy sampling''.

  \subsubsection{Outline.}
    %We start %in \autoref{part:prelims} 
      We begin in Part 1 by laying the foundations, with
      definitions for basic crypto concepts in \autoref{sec:basics}, and a
      primer on quantum computing concepts in \autoref{sec:primer}.
      In \autoref{sec:pseudocode}, we introduce a novel pseudocode notation,
      which will be helpful in reasoning about oracles and reductions.
      Then in \autoref{sec:rom}, we finally introduce our main objects of study,
      the quantum random oracles. Through illustrative examples we will be
      building intuition, aiming to get a ``feel'' for how they work.

    \hhnote{update the following once the v0.1 structure is settled}

    \autoref{part:classical} is all about the classical case. Here, we show
      several examples of security proofs in the random oracle model, all of
      which are conceptually simple and yield (relatively) tight security.
      The examples are lifted from the security proofs of two important
      transformations, the Fujisaki-Okamoto (FO) transformation (which transforms
      a weakly-secure public-key encryption scheme into a strongly-secure key
      encapsulation mechanism), and the Fiat-Shamir (FS) transformation (which
      transforms an interactive zero-knowledge protocol into a non-interactive
      signature scheme). 

    We meet these examples again in \autoref{part:tools}, where we use
      them to showcase the tools as we introduce them by showing that security
      holds also in the QROM, though usually with a security loss compared to the
      classical case, with the particular security loss depending on the tool used.
      \autoref{sec:o2h} is all about the One-way To Hiding (O2H) lemma, and its many variants: 
      how to use it to reprogram quantum random oracles, when to employ which
      version, and the impact of each version on the tightness of the proof. 
      Our running example here
      will be the second half of the FO transformation, known as the U-transform,
      and showing that if we start with a one-way secure (OW-CPA) deterministic PKE, we
      get an indistinguishable KEM secure against passive attacks (IND-CPA). 
      In \autoref{sec:extract}, we introduce the technique of
      Quantum Online-Extraction (QOE), and show how it can be used to prove
      that the U-transform \emph{also} upgrades to security against active attacks
      (IND-CCA).
      In \autoref{sec:adaptive}, we show how reprogramming can be made
      \emph{adaptive}, meaning that the reprogramming may now happen
      on-the-fly, on points that may (in full or in part) depend on
      adversarially-chosen input values. We use this to show
      how to go from a passively-secure Non-Interactive Zero-Knowledge (NIZK)
      protocol (such as constructed from a sigma protocol via the FS
      transformation) to an actively-secure signature scheme (EUF-CMA).
      Then, in \autoref{sec:compressed}, we delve into the details of 
      Compressed Random Oracles (CROs), and show how one might implement the
      technique to provide an efficient yet exact simulation of a quantum
      random oracle.
      This incredibly powerful technique takes some effort to wrap one's head
      around, and so we dedicate significant space to give gentle and
      intuitive explanations for all its moving parts, with several examples
      showing how queries are processed.
    \ifdraftversion
      Finally, in \autoref{sec:rewind}, we show how to lift rewinding arguments
      from the ROM to the QROM, with the central example being how the FS
      transformation transforms a kind of interactive zero-knowledge protocol known
      as a sigma protocol, to a NIZK.
      \hhnote{leave for version 2}
    \fi

    \ifdraftversion
    \autoref{part:proofs} is all about \emph{constructing the tools
      themselves}, which requires a deeper dive into the world of quantum
      information theory.
      \hhnote{and so on and so forth; leave for version 3}
    \fi

    The appendices contain extra material, including some more formal
      definitions in \autoref{sec:pseudocode}, and solutions to select
      exercises in \autoref{sec:solutions}.

    \hhnote{fix ``Sect.'' vs ``App.''}

  \clearpage

  \addcontentsline{toc}{part}{Part 1. Classical Security is Easy with Random Oracles}
  %\part{Setting the Stage}\label{part:prelims}
    \section{Cryptography Basics}\label{sec:basics}
  \todoenum[rough structure]{
    \item PKE and KEM: Definitions, OW vs IND, CPA vs CCA
    \item Signatures (and NIZKs?): Definition, EU-CMA
    \item Game hopping
  }

  \subsection{Sets, Lists, and Probabilities}
    Let us start by defining our conventions. 
      We will take the natural numbers $\NN$ to include $0$, and for $n \in \NN$, we
      will use $[n]$ to denote the set $\{0, \dots, n-1\}$. We will also use
      capital letters to denote powers of two, such as $N = 2^n$ (and $NM =
      2^{n+m}$).

    Functions are defined over bitstrings unless otherwise noted, and the set
      of all functions from the set of $m$-length bitsrings $\bin^m$ to the
      set of $n$-length $\bin^n$ bitstrings is denoted $\funcspace{m}{n}$.
      The $n$-by-$n$ identity matrix is denoted $\idmatrix[n]$, where we'll
      drop the subscript whenever dimension is clear from context.

    For a bitstring $x\in\{0,1\}^*$, we let $\card{x}$ denote its length;
      for a finite set $\Xspace$, we let $\card{\Xspace}$ denote its
      cardinality; and likewise for lists.
      For a list $\Xlist$, $\Xlist[i]$ retrieves its $i$th element,
      and two lists $\Xlist$ and $\Ylist$ may be combined by concatination to form
      a new list $\Zlist = \Xlist \concat \Ylist$. It will at times be useful
      to treat bitstrings as lists of binary values.

    Deterministic assignment to a (classical) variable is denoted $x \gets y$;
      uniform sampling from a set $\Xspace$ (and assignment of the result to
      $x$) is denoted $x \sample \Xspace$. Likewise, assignment of outputs from
      deterministic resp.\ probabilistic algorithms $X$, $Y$
      is denoted $x \gets X(\cdot)$ and $y \sample Y(\cdot)$, respectively.

    We use $\condprobsample{\mathsf{Code}}{\mathsf{Event}}{\mathsf{Condition}}$
      to denote the conditional probability of $\mathsf{Event}$ occurring
      when $\mathsf{Code}$ is executed, conditioned on $\mathsf{Condition}$
      (where the probability is taken over the random coins of
      $\mathsf{Code}$).
      We omit $\mathsf{Code}$ when it is clear from context and
      $\mathsf{Condition}$ when it is not needed. 

  \subsection{Algorithms and Oracles}
    Algorithms are assumed to be either classical~\cite{Turing36} or quantum Turing
      machines~\cite{Deutsch85} (although we will follow the convention of
      expressing quantum algorithms as quantum circuits~\cite{Yao93});
      which of the two we are talking about will usually be clear from context.
      An \emph{oracle} algorithm is an algorithm $A$ that additionally has black-box access to
      some oracle $\Oracle$, meaning it can only call the oracle on some input
      and receive its output; this is written $A^\Oracle$. 

    Unlike algorithms, oracles are typically not limited to bounded resources.
      Classical oracles take a classical value $x$ as input and return some
      value $y$, and can be deterministic (repeat queries return the same value) or
      probabilistic (the output depends on randomness internal to the oracle);
      in either case we write $y \gets \Oracle(x)$.
      Oracles can furthermore be either stateful (queries are remembered) or
      stateless. Clearly, stateless, deterministic oracles implement a function,
      and we may denote an oracle implementing the function $f$ as $\Oracle_f$.

    Quantum(-accessible) oracles\footnote{
      Some works distinguish between quantum-\emph{accessible} oracles,
      which evaluate a function on a superposition of inputs, and
      \emph{quantum} oracles, which apply some quantum operation on the input
      registers.
      We do not distinguish between these two kinds of oracle, and rather view
      the former as a special case of the latter.
      } are oracles that instead take a set of \emph{qubit registers} as input,
      apply some operation on them, before returning the registers to the
      querying party. The operation applied may be probabilistic in nature and
      may include measurements (which has the effect of transforming pure
      quantum states into mixed states\footnote{
        Such quantum operators are known as \emph{superoperators}
      }), though it will often be required to be unitary.
      For quantum oracle algorithms, we will take a classical oracle to be a
      quantum oracle that performs a computational-basis measurement on the
      registers before operating on them. Hence, we do not make a
      distinction between classical and quantum oracle access 
      in our notation, writing $A^\Oracle$ for either case. 

    In \autoref{subsec:pseudocode}, we define a quantum pseudocode that allows
      us to talk about general quantum oracles, which we apply to quantum random
      oracles in \autoref{sec:rom}. For now, we note that a quantum oracle
      $\Oracle_f$ that evaluates a function $f$ on superpositions of inputs are
      defined as the oracle that applies the unitary operation
      \[
        \ket{x,y} \to \ket{x,y \xor f(x)}\, ,
      \]
      where $\ket{x,y}$ are the computational basis states of the qubit
      input/output registers given to the oracle. (See \autoref{sec:primer} for
      an introduction to the basics of quantum computation.)

    %% MORE FROM SOA SOK
    %We use code-based experiments, where by convention all sets, lists,
      %and lazily sampled functions are initialized empty.
      %In our experiments, we will assume that no oracle calls are made (by the adversary)
      %with evidently out-of-bounds inputs (e.g.~list indices or key handles).

    %We use the shorthand $\Xset \setapp x$ for
      %the operation $\Xset \gets \Xset \cup \{ x \}$ and
      %$\Xlist \listapp x$ for appending the element $x$ to a list $\Xlist$.

    %We can make the
      %randomness $\rand$ of $Y$ explicit by writing
      %$x \gets Y(\, \cdot\, ; \rand)$, for previously sampled $\rand$.

    %We will also consider stateful algorithms, for instance when we write
    %  $\mssg \sample \mssgsamp[\samplestate](\dist)$
    %  the algorithm $\mssgsamp$ takes as state $\samplestate$ and
    %  as input the value $\dist$ and outputs $\mssg$,
    %  but it may change its state $\samplestate$ as part of its processing
    %  (code-snippet taken from \autoref{fig:IND-SOA}).

  \subsection{Public Key Encryption}
    First, whenever some object is generated
      dependent on some other public object, such as the public key $\pk$ being
      dependent on the public parameters $\params$ (which in turn depend on the
      security parameter $\secpar$), we assume without loss of generality that
      the new object includes descriptions of the prior (public) objects (so if you
      know $\pk$, then you also know $\params$). Similarly, we assume
      that if you know a secret key $\sk$ belonging to a keypair $(\pk, \sk)$, then 
      you also know $\pk$. This is w.l.o.g.\ in the sense
      that after a keypair $(\pk,\sk)$ is generated, one can always define the
      actual keypair to be $(\pk, (\pk, \sk))$; but it is not w.l.o.g.\ in
      the sense that this enfores that secret keys have unique public keys.\footnote{
        There exist PKE schemes that violate this, and
        there are even security proofs that rely on such non-unique key
        relations~\cite{EC:BelHofYil09}.
      }

    Adversaries are assumed to be generally quantum agents, having
      access to as much classical and quantum memory as they need, on which
      they can apply unitary operations and measurements as they
      please. (We will come upon cases later where it will be beneficial to restrict this
      assumption.) We will mostly think of these agents as being \emph{efficient} in the sense of
      running in time polynomial in the security parameter $\secpar$, although
      we will stick with a concrete formal treatment wherein asymptotics are
      avoided, and the running time of reductions are instead stated directly in terms
      of the running time of the original adversary. We may also
      consider \emph{unbounded} adversaries; these can be thought of as running
      in time (and space) exponential in the security parameter.

    If the adversary $\adv$ has access to a random oracle $\RO$, this is
      written $\adv^{\RO}$, and likewise for any other oracles it may have
      access to. (Our notation does not distinguish between quantum and
      classical oracles.) Similarly, if a reduction $\bdv$ has black-box access
      to an adversary $\adv$ (meaning it may control its input and read its output, but not
      otherwise read or affect its internal workings), this is written $\bdv^\adv$.

    \subsubsection{Defining security.}
      The standard security notion for public-key encryption is usually taken to
        be ``Indistinguishability under Chosen-Ciphertext Attacks'', or $\indcca$,
        but for the sake of this tutorial we will stick to the weaker 
        ``Indistinguishability under Chosen-\emph{Plaintext} Attacks'' ($\indcpa$). Since
        this is a public-key encryption scheme, the chosen-plaintext attacks can
        be performed ``offline'' by the adversary simply by running the
        encryption algorithm itself using the provided public key. 
        %Therefore, unlike in the symmetric encryption case, we do not need to provide 
        %the adversary with an honest encryption via an oracle.
        % Eh, it gets this anyway, by letting m0 = m1, making the point moot.

      On the left of \autoref{fig:ind-ow} is the experiment
        defining our (single-user) $\indcpa$ security notion. The game starts by sampling
        parameters $\params$, a keypair
        $(\pk,\sk)$, and a challenge bit $\chalbit$, before handing the public
        key to the adversary $\adv$. If the construction in question uses a
        random oracle $\RO$, then $\adv$ also gets oracle access to $\RO$. 

      At some point, the adversary decides on two equal-length messages
        $\mssg_0$ and $\mssg_1$, and call its \emph{challenge encryption} oracle $\Oenc$
        (\autoref{fig:ind-ow}, middle). This oracle then provides an honest
        encryption of the message $\mssg_{\chalbit}$ and returns a challenge
        ciphertext $\chalctxt$. If the adversary ever calls
        the challenge oracle on messages of differing lengths, then it returns the
        special ``fail'' symbol $\fail$ instead; this is due to the
        assertion that a scheme should be considered
        secure even if ciphertexts leak information about message lengths.
        In fact, you may notice from \autoref{fig:bre} that
        our BRE scheme leaks the message length exactly, since $\demctxt$
        must necessarily have the same length as $\mssg$. 
        
      We are in the \emph{multi-challenge} setting, and so the adversary may
        continue to make as many calls to $\Oenc$ as it likes,
        in addition to arbitrary calls to the random oracle $\RO$. 
        (If the adversary is restricted to a single call to $\Oenc$, we call this
        the \emph{single-challenge setting}.)
        Once it is satisfied, $\adv$ halts with a guess $\advbit \in
        \bin$, and wins if it guessed $\chalbit$ correctly.

      We measure security in terms of $\adv$'s \emph{distinguishing advantage},
        i.e.\ how much better it can do than random guessing.

      \begin{figure}[t]
        \centering
        \begin{pchstack}[center]
          \procedure{$\experiment{\indcpa}{\pkescheme}$}{
            \params \sample \pkpm(\secparam) \\
            (\pk, \sk) \sample \pkkg(\params) \\
            \chalbit \sample \bin \\
            \advbit \sample \adv^{\Oenc(,\RO)}(\pk) \\
            \pcreturn \advbit = \chalbit
          }
    
          \pchspace
    
          \procedure{$\indOenc(\mssg_0, \mssg_1) \pccomment{\ind}$}{
            % Single-challenge:
              %\pcif \chalctxt \neq \emptystring \lor \abs{\mssg_0} \neq \abs{\mssg_1} \pcthen \\ 
            \pcif \abs{\mssg_0} \neq \abs{\mssg_1} \pcthen \pcreturn \fail \\
            \chalctxt \sample \pkenc_\pk(\mssg_\chalbit) \\
            \pcreturn \chalctxt
          }
    
          \pchspace
    
          \procedure{$\experiment{\owcpa}{\kemscheme}[(\adv)]$}{
            \params \sample \kempm(\secparam) \\
            (\pk, \sk) \sample \kemkg(\params) \\
            i \gets 0 \\
            (i,\advdemkey_i) \sample \adv^{\Oenc(,\RO)}(\pk) \\
            \pcreturn (\advdemkey_i = \chalkeyi)
          }
    
          \pchspace
    
          \procedure{$\owOenc \pccomment{\oneway}$}{
            i \gets i+1 \\
            (\chalkeyi, \chalctxt) \sample \kemenc_\pk \\
            \pcreturn \chalctxt
          }
        \end{pchstack}
      \caption{
        \label{fig:ind-ow}
        The experiments defining (multi-challenge) $\indcpa$ for PKE (left), and
        $\owcpa$ for KEM (right). If a random oracle $\RO$ is present, then the
        adversary gets access to this in addition to the challenge oracle $\Oenc$.
      }
      \end{figure}

      \begin{definition}[Distinguishing Advantage]
        The distinguishing advantage of an adversary $\adv$ against a
          public-key encryption scheme $\pkescheme$, under chosen ciphertext
          attacks, is defined as
        \begin{align*}
          \advantage{\indcpa}{\pkescheme} &\coloneqq \abs{
            2 \cdot \prob{\experiment{\indcpa}{\pkescheme}} - 1
          }\, ,
          \intertext{or, equivalently,}
          \advantage{\indcpa}{\pkescheme} &\coloneqq \abs{
            \prob{\experiment{\indcpa}{\pkescheme}} 
            - \prob{\neg\experiment{\indcpa}{\pkescheme}}
          }\, ,
          \intertext{or, equivalently,}
          \advantage{\indcpa}{\pkescheme} &\coloneqq \abs{
            \condprobsample{
              \experiment{\indcpa}{\pkescheme}
            }{0 \sample \adv}{\chalbit = 0}
            - \condprobsample{
              \experiment{\indcpa}{\pkescheme}
            }{0 \sample \adv}{\chalbit = 1}
          }\, .
        \end{align*}
        where $\experiment{\indcpa}{\pkescheme}[]$ is as given in
        \autoref{fig:ind-ow}.
      \end{definition}

      \noindent
      As is easiest seen from the first of the above three equations, the
        advantage is just the probability that $\adv$ wins the game, minus the
        probability of winning based on a random guess, scaled so that the
        advantage is a value defined on the interval from $0$ to $1$,
        where $0$ means $\adv$ has no advantage (e.g.\ if the messages are perfectly hidden), 
        and $1$ means that $\adv$ is a perfect distinguisher (the scheme is
        completely broken, or $\adv$ is unbounded).

  \subsection{Key Encapsulation Mechanisms}

  \subsection{Digital Signature Algorithms}

  \subsection{Security Reductions}
    \subsubsection{Kinds of reductions.}
      We say that a reduction is \emph{simple} if it only runs the adversary once,
        and \emph{fully black-box} if it runs the adversary from start to finish,
        answering oracle queries along the way, and otherwise only providing its
        input and reading its output; if, as is typically the case, the reduction
        is both simple and fully black-box, we call it an SFBB reduction. 

      We will later see some reductions that \emph{rewind} the adversary, and run
        it again from a point chosen by the reduction. We model this as the usual
        black-box access with an additional rewinding \emph{interface}, with which
        it can essentially tell the adversary to rewind itself to a specified point
        in its computation; we refer to it as a \emph{rewinding black-box} (RBB)
        reduction.


    \subsubsection{Game Hopping.}
      %\todoenum[Provide explanations of:]{
      %  \item Basic background on game hopping and advantages, stuff like
      %    difference vs.\ times two minus 1, with or
      %    without absolute values, the fundamental lemma of game hopping.
      %  \item Maybe write something about asymptotic vs.\ concrete security.
      %}
      For our security proofs, we will be working in the game-hopping
        paradigm of Shoup~\cite{EPRINT:Shoup04}.
        The fundamental lemma of game hopping states that the distinguishing advantage
        of an adversary can be bounded in terms of their distinguishing advantage
        over several intermediary game ``hops''. 

      \begin{lemma}[Fundamental Lemma of Game Hopping]\label{lemma:game-hopping}
        For any adversary $\adv$ tasked with distinguishing $\game{0}$ from
        $\game{n}$, and for any games $\game{1}, \ldots \game{n-1}$,
        there exist adversaries $(\adv_1, \ldots, \adv_n)$ such that
        \[ 
          \abs{
            \probsample{\game{0}(\adv)}{1 \sample \adv} - 
            \probsample{\game{n}(\adv)}{1 \sample \adv}
          } \leq \sum_{i = 1}^n \abs{
            \probsample{\game{i-1}(\adv)}{1 \sample \adv_i} - 
            \probsample{\game{i}(\adv)}{1 \sample \adv_i}
          }\, .
        \]
      \end{lemma}

      \noindent
      The lemma states that, if $\adv$ were able to distinguish, then
        at least one of $\adv_1, \ldots, \adv_n$ must be able to
        distinguish too---and vice versa, that if none of the $\adv_i$ are able to
        distinguish, then neither is $\adv$.
        (In fact, the $\adv_i$ that achieve this will in essence be just copies
        of $\adv$.)

      The proof is remarkably simple, reproduced here for those who haven't 
        seen it.

      \begin{proof}
        For all $i$ from $1$ to $n-1$, insert $0 = \probsample{\game{i}(\adv_i)}{1 \sample \adv_i} -
        \probsample{\game{i}(\adv_i)}{1 \sample \adv_i}$ into the absolute value
        on the left hand side, then apply the union bound.
        \qed
      \end{proof}

    %\paragraph{Quantum notation.}
    %  %  While we only assume knowledge of basic linear algebra, this tutorial 
    %  %    still concerns itself with several fundamentally quantum concepts. To
    %  %    help bridge this divide, we provide a self-contained primer on quantum
    %  %    computing theory in
    %  %    \autoref{sec:primer}, with explanations and definitions for all
    %  %    the quantum concepts and conventions that we will need. A summary of our
    %  %    notation is given in \autoref{tab:notation}.
    %  Precise definitions for our quantum-notational conventions are provided in
    %    \autoref{sec:primer},
    %    along with a gentle introduction to the relevant quantum concepts,
    %    summarized in \autoref{tab:notation}.

      \hhnote{
        Add: $\pcparse$ keyword. (Allow $\pcif \pcparse$ too? By having $\pcparse$
        return a boolean value? Hmm, tricky.) And: $\red\bad$ abort convention.
      }
      \hhnote{
        Add: Adversaries are always assumed to halt in due time, even when the
        environment fails to match their expectation.
      }

    \subsubsection{Bad Events.}
      \hhnote{
        Add identical-until-bad lemma?
      }

    \clearpage
    % Classical lazy-sampling
\begin{figure}[t]
  \centering
  \begin{pchstack}[center]
    \procedure{$\RO(x)$}{
      \pcif \Db[x] = \bot \pcdo \Db[x] \sample \bin^m \\
      \pcreturn \Db[x]
    }
  \end{pchstack}
  \caption{
    \label{fig:lazy-sampling}
    A lazy-sampled random oracle from $n$ to $m$ bits. 
    Here, $\Db$ is a query database. Note that the oracle is \emph{stateful},
    since the database is maintained between calls.
  }
\end{figure}
%
\section{The Random Oracle Model}\label{sec:rom}
  Random oracles are nice. They are also very easy to simulate, a process known as
    lazy-sampling, as shown in \autoref{fig:lazy-sampling}.

  \subsection{From a One-way PKE to an Indistinguishable KEM: The $\U$-transform}
    The Fujisaki-Okamoto transformation~\cite{C:FujOka99}, which underlies many post-quantum
      KEMs including the now-standardized ML-KEM (Kyber)~\cite{Kyber}, can be
      framed as the composition of two 
      transformations, \T and \U, as defined by Hofheinz
      et al.~\cite{TCC:HofHovKil17} (hereafter HHK). 
      Starting from a public-key
      encryption scheme, the first transformation (\T) derandomizes the
      PKE by having its randomness be derived from a hash of the message (and possibly
      the public key), and additionally adds a re-encryption check to the
      decryption algorithm, while the second transformation builds a KEM from a
      deterministic PKE by sampling a uniform message, encrypting it, and
      hashing it to produce the ephemeral key $\eK$.
  
    In the following, we will focus on this second transformation, $\U$,
      and build our KEM starting with a deterministing PKE that satisfies
      one-way security.\footnote{
        HHK defined several variants of the \U-transform that differ from the variant discussed here
    		in how they handle invalid ciphertexts and in what they hash into the
        output key; our variant corresponds to \UExpMess in their notation.
        %\hhnote{
        %  Should we do the more general \UExpMess transform instead? Or does that
        %  just muddy the message?
        %}
      } (See App.~\ref{sec:ttransform} for details on the \T-transform.)

    \hhnote{
      Leave this version without the extra re-encryption check (which also
      matches the original \U), and define a new \U with the re-encryption check
      when we later want CCA security.
    }
   
    \begin{definition}[\U-transform]
      To a deterministic PKE scheme $\dPKE = (\KG, \Enc, \Dec)$ and a hash
      function $\ROkey: \MSpace \to \KSpace$,
      we associate the KEM $\KEM \coloneqq \U[\PKE, \ROkey]$ defined in \autoref{fig:def:U-transform}.
      \begin{figure}[ht]
   	\begin{center}
   		\fbox{
   			\small
   			\nicoresetlinenr
   			%\begin{minipage}[t]{3.3cm}
   			%	\underline{$\UKG()$}
   			%	\begin{nicodemus}
   			%		\item $(\pk, \sk) \sample \KG()$
   			%		%\item $\rejectionSeed \sample \MSpace$
   			%		%\item $\sk' \leftarrow (\sk, \rejectionSeed)$
   			%		\item \pcreturn $(\pk, \sk)$
   			%	\end{nicodemus}
   			%\end{minipage}
   			%
   			%\quad
   			
   			\begin{minipage}[t]{3cm}
           \underline{$\UEncaps_\pk()$}
   				\begin{nicodemus}
   					\item $m \sample \MSpace$
   					\item $c \gets \Enc(\pk, m)$
   					\item $\eK \gets \ROkey(m)$
   					\item \pcreturn $(\eK, c)$
   				\end{nicodemus}
   			\end{minipage}
   			
   			\quad
   			
   			\begin{minipage}[t]{3.6cm}
          \underline{$\UDecaps_\sk(c)$}
   				\begin{nicodemus}
   					%\item Parse $\sk' := (\sk, \rejectionSeed)$
            \item $m' \gets \Dec_\sk(c)$
   					\item \pcif $m' = \bot \pcthen \pcreturn \bot$
   						%\item \pcreturn $\ROkey(\rejectionSeed, c)$
   					\item \pcreturn $\ROkey(m')$
   				\end{nicodemus}
   			\end{minipage}
   		}
   	\end{center}
   	\caption{
   	  \label{fig:def:U-transform}			
        Algorithms of $\U[\PKE, \ROkey] \eqqcolon (\KG, \UEncaps, \UDecaps)$.
     }
   \end{figure} 
   \end{definition}
  
  \subsection{What Goes Wrong}
    \todoenum[rough structure]{
      \item A discussion about why the previous proofs do not yield
        post-quantum security 
      \item Where the proofs break down, and why they do not admit direct
        lifting (not history-free)
    }

    \clearpage
    \section{\titleUtransformRom}\label{sec:utransform}
  
  We will now meet the Fujisaki--Okamoto (FO) transformation~\cite{C:FujOka99},
  	a PKE-to-KEM recipe that became quite popular in the post-quantum-context --
  	it is used as part of many post-quantum KEMs,
  	including the now-standardized ML-KEM (Kyber)~\cite{FIPS203}. 
  
  We first focus on the part of FO that turns a deterministic PKE scheme into a KEM.
	  Following the literature~\cite{TCC:HofHovKil17}, this part of FO will be called the \U-transform.
	  %, as defined by Hofheinz et al. (hereafter HHK).
	  To have a simple ROM example, we will -- for now -- only prove that the obtained KEM
	  achieves key indistinguishability against \emph{passive} adversaries.
	  This proof will already teach us several useful tricks you can play in the ROM.
	  These tricks foreshadow and motivate the QROM tools that will be introduced in~\cref{sec:o2h}
	  \todo{\cref{sec:o2h} vs \cref{sec:C-O2H}}.

  At a later point\todo{add pointer to where},
    we will also prove key indistinguishability against \emph{active} adversaries.
    There also exists a second part of FO that turns probabilistic PKE schemes
    into ones that fit the \U-transform's requirements.
    (This part is called the \T-transform, in case you're curious, see App.~\ref{sec:ttransform}.)
%    (\T) de-randomizes it by first deriving the random coins by hashing
%    the message (and possibly the public key), then running the encryption
%    algorithm using those coins; then, after decrypting, it re-derives the
%    random coins and re-encrypts the message to verify that resulting
%    ciphertext matches the input. 

  In short, the \U-transformation builds a KEM from a PKE scheme by sampling a random message,
    encrypting it, and hashing it (possibly together with the public key and ciphertext) to derive a
    symmetric key $\eK$. This is the transformation that will serve as our
    running example.\footnote{
      When Hofheinz et al.~\cite{TCC:HofHovKil17} first defined
      the \T- and \U-transforms, they actually defined several \U-variants;
      %that differ from the variant discussed here
  		%in how they handle invalid ciphertexts and in what they hash into the
      %output key; 
      \cref{def:U-transform} corresponds to
      \UExpMess in their notation (message-only hashes, explicit decryption failures).
    }

  \begin{definition}[\U-transform]\label{def:U-transform}
    Let $\dPKE = (\KG, \Enc, \Dec)$ be a deterministic public-key encryption scheme, and let 
    $\ROkey: \bin^* \to \bin^m$ be a hash function.
    Then the \U-transform is the key-encapsulation mechanism 
    $\U[\PKE, \ROkey] \eqqcolon (\KG, \UEncaps, \UDecaps)$, as defined in \cref{fig:def:U-transform}.
    \begin{figure}[ht]
   	  \begin{center}
  		  \fbox{
  		  	\small
  		  	\nicoresetlinenr
  		  	%\begin{minipage}[t]{3.3cm}
  		  	%	\underline{$\UKG()$}
  		  	%	\begin{nicodemus}
  		  	%		\item $(\pk, \sk) \sample \KG()$
  		  	%		%\item $\rejectionSeed \sample \MSpace$
  		  	%		%\item $\sk' \leftarrow (\sk, \rejectionSeed)$
  		  	%		\item \pcreturn $(\pk, \sk)$
  		  	%	\end{nicodemus}
  		  	%\end{minipage}
  		  	%
  		  	%\quad
  		  	
  		  	\begin{minipage}[t]{3cm}
            \underline{$\UEncaps_\pk()$}
  		  		\begin{nicodemus}
              \smallskip
  		  			\item $x \sample \bin^n$
              \smallskip
  		  			\item $c \gets \Enc(\pk, x)$
              \smallskip
  		  			\item $\eK \gets \ROkey(x)$
              \smallskip
  		  			\item \pcreturn $(\eK, c)$
  		  		\end{nicodemus}
  		  	\end{minipage}
  		  	
  		  	\quad
  		  	
  		  	\begin{minipage}[t]{2.5cm}
           \underline{$\UDecaps_\sk(c)$}
  		  		\begin{nicodemus}
              \smallskip
  		  			%\item Parse $\sk' := (\sk, \rejectionSeed)$
              \item $x' \gets \Dec_\sk(c)$
              \smallskip
  		  			\item \pcif $x' = \bot :$ 
              \smallskip
              \item   \quad \pcreturn $\bot$
              \smallskip
  		  				%\item \pcreturn $\ROkey(\rejectionSeed, c)$
  		  			\item \pcreturn $\ROkey(x')$
  		  		\end{nicodemus}
  		  	\end{minipage}
  		  }
  	  \end{center}
  	  \caption{
  	    \label{fig:def:U-transform}			
         Algorithms of $\U[\PKE, \ROkey]$.
      }
    \end{figure} 
  \end{definition}

  %\hhnote{
  %  Leave this version without the extra re-encryption check (which also
  %  matches the original \U), and define a new \U with the re-encryption check
  %  when we later want CCA security.
  %}

  As explained in \cref{sec:basics}, One-way security under Chosen
    Plaintext Attacks, or \owcpa, is really a minimal security
    requirement for PKE schemes, and should be considered too weak a security
    goal for most use-cases. 
    %The goal of One-Wayness (OW) simply says that
    %given a ciphertext, it should be hard to recover the plaintext, and the
    %power of Chosen Plaintext Attacks (CPA), the ability to encrypt any
    %message under the key you are targetting, are of course ever-present when
    %encryption is public.
    Much preferred would be the stronger security requirement of
    \emph{indistinguishability} under chosen-plaintext attacks, or
    \indcpa.\footnote{
      The standard security requirement for deployable KEMs is an even stronger one,
      namely indistinguishability under chosen-\emph{ciphertext} attacks, or
      \indcca. 
      As it turns out, the \U-transform also achieves this, provided the
      re-encryption check from the \T-transform is carried forward; this will
      be covered in later chapters (to appear).
    } 

  However, stronger security assumptions usually mean less efficient
    schemes, as the parameters must be increased to compensate for the
    increased attack surface. What's worse, we are taking as our starting
    point a \emph{deterministic} encryption scheme. Such a PKE scheme could
    \emph{never} achieve indistinguishability, since the adversary can simply
    re-encrypt the two candidate messages, and check which of them
    encrypts to the challenge ciphertext.

  And yet, taking our deterministic PKE scheme through the \U-transform
    does in fact give us a KEM that achieves indistinguishability
    \ldots and the object that allows us to cross that chasm
    %between one-wayness and indistinguishability 
    is precisely the \emph{random oracle}.

  The proof strategy will rely on the fact that our reduction can
    fully control the random oracle for the adversary, inspecting the queries
    before answering them via the lazy-sampling procedure from \cref{sec:lazy-sampling}. 
    Intuitively speaking, in order to successfully distinguish, the adversary
    will have to have called the random oracle on one of
    the messages used to derive the ephemeral key. 
    The reduction can recognize this the moment it happens, by re-encrypting
    each message and checking if the encryption matches the challenge it was
    issued; if so, it wins at the one-wayness game by outputting that message.

    %able to recognize the correct solution; this fact allows us to get
    %a tight bound in this case. (Otherwise we would have to guess which of the
    %calls to the random oracle contained the solution we are looking for,
    %which would incur a tightness loss related to the probability that
    %we guess correctly.)

  To summarize, if $\dPKE$ satisfies one-wayness
    (\cref{def:owcpa}), then $\KEM$ satisfies indistinguishability
    (\cref{def:indcca}), as captured in the following theorem.

  \begin{theorem}[\dPKE is \owcpa $\To$ \KEM is \indcpa]
    \label{thm:cparom}
    Let \dPKE be a deterministic PKE scheme, let \RO be a random oracle, 
    and let the $\KEM = \U[\dPKE, \RO]$.
    Then there is a simple, fully black-box reduction $\bdv$, such that for
    every adversary $\adv$,
    %that makes at most $\qe$ calls to the challenge oracle, 
    \[
      \advantage{\indcpa}{\KEM}[\left(\adv^{\RO}\right)] 
        \leq \advantage{\owcpa}{\dPKE}[\left(\bdv^{\adv}\right)]\, .
    \]
    The adversary $\bdv$ runs in the same time as $\adv$, plus the %(relatively small) 
    overhead of lazy-sampling a random oracle.
  \end{theorem}

  \begin{proof} %[\cref{thm:cparom}]
    We want to construct a $\bdv$ that exploits any winning
      strategy of $\adv$ in such a way that if $\adv$ \emph{would} win at the
      indistinguishability game, then $\bdv$ \emph{actually} wins
      at the one-wayness game.
      However, the constructed adversary $\bdv$ will not be able simulate the
      real ($b=0$) world of the KEM indistinguishability game, since it does
      not know which message the challenge ciphertext encrypts -- that's what
      it is trying to figure out! 
      %hence, it can not hash the message to produce the
      %correct symmetric key for $\adv$.
      It can, however, easily simulate the random ($b=1$) world, since here the
      symmetric key given to the adversary is simply uniformly sampled.

    But in order to argue that $\adv$ can not distinguish $\bdv$'s simulation
      from the real-world game, we need to alter the game in such a way that the
      move from real to random keys is completely hidden from the view of $\adv$.
      We do this by inserting two intermediary games between the real-world and the
      random-world games, as illustrated in \cref{fig:cparom-hops-illustrated}.
      These two middle games formalize the intuition that, thanks to the \U-transform,
      the adversary will have to call the random oracle on the correct message
      in order to notice that it is not in the real world: If it
      does, then a boolean flag $\bad$ is set to true. Assuming the $\bad$ flag is
      never set to true, the game proceeds exactly as the real-world game, and
      thanks to \cref{lemma:bad}, we know that the noticeable difference
      between these games is bounded exactly by the probability that $\bad$ is
      set to true. Then, \emph{conditioned} on $\bad$ remaining false,
      we can exchange the real symmetric key for a randomly
      sampled one, and be sure that the adversary will never be able to
      notice.
      %Therefore, the second game hop from real to random keys is ``for free'',
      %since the two games are equal in terms of probabilities. 


      \crakzfigcode{cparom-hops-illustrated}{
        Proof outline for \cref{thm:cparom}. The arrows show the changes
        incurred by each game hop, and the braces show the difference in
        win probabilities between each game.
        %Here, $\indcpa_0$ resp.\ $\indcpa_1$ is shorthand for
        %the game $\experiment{\indcpa}{\pkescheme}[]$ (\cref{fig:ind-ow}) instantiated with challenge bit
        %$\chalbit = 0$ (left world) resp.\ $\chalbit = 1$ (right world).
      }
  
    The four games are given in pseudocode in \cref{fig:cparom-hops}, in a
      layout matching \cref{fig:cparom-hops-illustrated}, and with
      changes between games highlighted.
      To make the argument formal, start by recalling the definition of a
      distinguishing advantage; we will use the ``difference'' version of
      distinguishing advantages, as
      it fits more cleanly with this ``game hopping'' style of proof.
      (For the sake of simplicity, we suppress the superscript dependencies
      $\RO$ and $\adv$ in the following.)
      \[
      	\advantage{\indcpa}{\KEM} \coloneqq \abs{
      		\condprobsample{
      			\experiment{\indcpa}{\KEM}
      		}{1 \gets \adv}{\chalbit = 0}
      		- \condprobsample{
      			\experiment{\indcpa}{\KEM}
      		}{1 \gets \adv}{\chalbit = 1}
      	}\, .
      \]
      The game $\game 1$ is just the real-world ($\chalbit=0$) $\indcpa$ game, with the
      construction of \KEM written out explicitly, together with a lazy-sampled
      random oracle. Similarly, $\game 4$ is just the random-world
      ($\chalbit=1$) $\indcpa$ game with the \KEM construction explicit. We
      can therefore write
      \[
      	\advantage{\indcpa}{\KEM} = \abs{
      		\probsample{ \game{1}(\adv) }{1 \gets \adv}
      		- \probsample{ \game{4}(\adv) }{1 \gets \adv}
      	}\, .
      \]
      We are of course always free to ``add zero'' by adding and subtracting
      the same term, so this is equal to:
      \begin{align*}
        \advantage{\indcpa}{\KEM} &= 
          |\probsample{ \game{1}(\adv) }{1 \gets \adv} 
            - \probsample{ \game{2}(\adv) }{1 \gets \adv} \\
          &\qquad + \probsample{ \game{2}(\adv) }{1 \gets \adv} 
            - \probsample{ \game{3}(\adv) }{1 \gets \adv} \\
          &\qquad + \probsample{ \game{3}(\adv) }{1 \gets \adv} 
            - \probsample{ \game{4}(\adv) }{1 \gets \adv}|\, .
      \end{align*}
      Applying the triangle inequality then gives us an upper-bound 
        on the original distinguishing advantage as the sum of the distinguishing
        advantages between successive games:
      \begin{align*}
        \advantage{\indcpa}{\KEM} &\leq 
          \abs{\probsample{ \game{1}(\adv) }{1 \gets \adv} 
            - \probsample{ \game{2}(\adv) }{1 \gets \adv}} \\
          &\qquad + \abs{\probsample{ \game{2}(\adv) }{1 \gets \adv} 
            - \probsample{ \game{3}(\adv) }{1 \gets \adv}} \\
          &\qquad + \abs{\probsample{ \game{3}(\adv) }{1 \gets \adv}
            - \probsample{ \game{4}(\adv) }{1 \gets \adv}}\, .
      \end{align*}

      Let us analyze each distinguishing advantage in turn.
        First, instead of deriving the key from the hash of the message as
        usual, $\game 2$ instead samples $\eK^*$ uniformly at random at the start
        of the game and gives it to $\adv$. Then, if $\RO$ is ever called with
        the encrypted message $x^*$, which \emph{should} have been hashed to
        derive $\eK^*$ at the start of the game, the oracle database is updated such that
        $\Db[x^*] = \eK^*$. Since the random oracle was lazy-sampled anyway,
        there is no way for $\adv$ to tell if the internal oracle database was
        updated right away, or only when it is needed. From the point of view
        of $\adv$, these games therefore behave identically, and
      \begin{align*}
        \abs{\probsample{ \game{1}(\adv) }{1 \gets \adv} 
          - \probsample{ \game{2}(\adv) }{1 \gets \adv}} = 0\, . 
      \end{align*}
      Two more things to notice about $\game 2$: First, the oracle interface
        could have simply done the above check by checking whether the input $x$
        matches the sampled message $x^*$; instead, it takes the roundabout
        route of re-encrypting the input and checking whether the resulting
        ciphertext is equal to $\ctxt^*$. Since \dPKE is deterministic
        (and by assumption perfectly correct), these checks lead to identical
        outcomes; but of course, only the latter check will be available to $\bdv$,
        who does not know $x^*$. Second, in addition to the ``delayed
        programming'' of the random oracle, the game introduces the $\bad$
        flag in line $10$. Initialized to $\false$, this flag is only set to
        $\true$ in the case that the database is actually updated to hold the
        sampled $\eK^*$, or in other words, if the random oracle is called on
        the encrypted message $x^*$. Other than this, the flag has no effect on the
        behaviour of the game.

      Compared to the changes introduced in $\game 2$, the change from 
        $\game 2$ to $\game 3$ is deceptively simple: When the random oracle is
        called on $x^*$, after $\bad$ is set to $\true$, do \ldots nothing.
        This is where $\game 2$ would program the random oracle with $\eK^*$,
        thus ensuring that the random oracle remained consistent with the
        real world. Now we are removing this programming, which means that
        once $\bad$ is set to true,
        the game is \emph{noticably} inconsistent with the real-world game:
        If $\adv$ was able to recover $x^*$ from $\ctxt^*$, then they will
        quickly notice that $\RO(x^*)$ does not evaluate to $\eK^*$ (except
        if, by freak accident, the random value drawn by the lazy-sampling
        procedure happens to be exactly $\eK^*$).

      Observe how $\game 2$ and $\game 3$ are by design exactly identicaly
        up until the point that $\bad$ is set to $\true$. This means that we
        can use \cref{lemma:bad} to bound their difference as the probability
        that $\bad$ is set to $\true$ at the end of either game:
      \begin{align*}
        \abs{\probsample{ \game{1}(\adv) }{1 \gets \adv} 
          - \probsample{ \game{2}(\adv) }{1 \gets \adv}} \leq
          \probsample{\game{2}(\adv)}{\bad} = \probsample{\game{3}(\adv)}{\bad}\, . 
      \end{align*}

    %%% Game hops, single-challenge version
    \begin{figure}[t]
    	\begin{center}
    	\resizebox{\textwidth}{!}{%
    	\begin{tabular}{@{}c@{\hspace{0em}}c@{}}
        % TOP LEFT: Game 1
    		\fbox{
    			\small
    			\nicoresetlinenr
          \begin{minipage}[t][3.4cm]{5cm}	
            \underline{\gameHeading $\game 1(\adv)$}
            \smallskip
    				\begin{nicodemus}
              \item $(\pk, \sk) \gets \KG()$
              \smallskip
              \item $x^* \sample \bin^n$
              \smallskip
              \item $\ctxt^* \assign \Enc_{\pk}(x^*)$
              \smallskip
              \item $\eK^* \gets \RO(x^*)$ 
              \smallskip
              \item $\advbit \gets \adv^{\RO}(\pk, \eK^*, \ctxt^*)$ 
              \smallskip
              \item $\pcreturn \advbit$
    				\end{nicodemus}
    			\end{minipage}
    
          \begin{minipage}[t]{3.6cm}
    				\underline{\oracleHeading $\RO(x)$}
            \smallskip
    				\begin{nicodemus}
              \item $\pcif \Db[x] = \bot :$ 
              \smallskip
              \item   \quad $\Db[x] \sample \bin^m$
              \smallskip
              \item $\pcreturn \Db[x]$
    				\end{nicodemus}	
    			\end{minipage}
    		}
        &
        %% TOP RIGHT: Game 4
    		\fbox{
    			\small
    			\nicoresetlinenr
          \begin{minipage}[t][3.4cm]{4.2cm}	
            \underline{\gameHeading $\game 4(\adv)$}
            \smallskip
    				\begin{nicodemus}
              \item $(\pk, \sk) \gets \KG()$
              \smallskip
              \item $x^* \sample \bin^n$
              \smallskip
              \item $\ctxt^* \assign \Enc_{\pk}(x^*)$
              \smallskip
              \item \dbox{$\eK^* \sample \bin^m$}
              \smallskip
              \item $\advbit \gets \adv^{\RO}(\pk, \eK^*, \ctxt^*)$ 
              \smallskip
              \item $\pcoutput \advbit$
    				\end{nicodemus}
    			\end{minipage}
    
    			\begin{minipage}[t]{3.6cm}
    				\underline{\oracleHeading $\RO(x)$}
            \smallskip
    				\begin{nicodemus}
              \item \sout{$\ctxt \assign \Enc_\pk(x)$}
              \smallskip
              \item \sout{$\pcif \ctxt = \ctxt^* :$}
              \item \quad \sout{$\red{\bad \assign \true}$}
              \smallskip
              \item $\pcif \Db[x] = \bot :$ 
              \smallskip
              \item   \quad $\Db[x] \sample \bin^m$
              \smallskip
              \item $\pcreturn \Db[x]$
    				\end{nicodemus}	
    			\end{minipage}
    		}
        \\[2ex]
        % BOTTOM LEFT: Game 2
    		\fbox{
    			\small
    			\nicoresetlinenr
          \begin{minipage}[t][4.2cm]{5cm}	
            \underline{\gameHeading $\game 2(\adv)$}
            \smallskip
    				\begin{nicodemus}
              \item \fbox{$\bad \assign \false$}
              \smallskip
              \item $(\pk, \sk) \gets \KG()$
              \smallskip
              \item $x^* \sample \bin^n$
              \smallskip
              \item $\ctxt^* \assign \Enc_{\pk}(x^*)$
              \smallskip
              \item \sout{$\eK^* \gets \RO(x^*)$} \fbox{$\eK^* \sample \bin^m$}
              \smallskip
              \item $\advbit \gets \adv^{\RO}(\pk, \eK^*, \ctxt^*)$ 
              \smallskip
              \item $\pcoutput \advbit$
              \smallskip
    				\end{nicodemus}
    			\end{minipage}
    
    			\begin{minipage}[t]{3.6cm}
    				\underline{\oracleHeading $\RO(x)$}
            \smallskip
    				\begin{nicodemus}
              \item \fbox{$\ctxt \assign \Enc_\pk(x)$}
              \item \fbox{$\pcif \ctxt = \ctxt^* :$} 
              \item \quad \fbox{$\red{\bad \assign \true}$} 
              \item \quad \fbox{$\Db[x] \assign \eK^*$} 
              \smallskip
              \item $\pcif \Db[x] = \bot :$ 
              \smallskip
              \item   \quad $\Db[x] \sample \bin^m$
              \smallskip
              \item $\pcreturn \Db[x]$
    				\end{nicodemus}	
    			\end{minipage}
    		}
        &
        % BOTTOM RIGHT: Game 3
    		\fbox{
    			\small
    			\nicoresetlinenr
          \begin{minipage}[t][4.2cm]{4.2cm}	
            \underline{\gameHeading $\game 3(\adv)$}
            \smallskip
    				\begin{nicodemus}
              \item $\bad \assign \false$
              \smallskip
              \item $(\pk, \sk) \gets \KG()$
              \smallskip
              \item $x^* \sample \bin^n$
              \smallskip
              \item $\ctxt^* \assign \Enc_{\pk}(x^*)$
              \smallskip
              \item $\eK^* \sample \bin^m$
              \smallskip
              \item $\advbit \gets \adv^{\RO}(\pk, \eK^*, \ctxt^*)$ 
              \smallskip
              \item $\pcoutput \advbit$
    				\end{nicodemus}
    			\end{minipage}
    
    			\begin{minipage}[t]{3.6cm}
    				\underline{\oracleHeading $\RO(x)$}
            \smallskip
    				\begin{nicodemus}
              \item $\ctxt \assign \Enc_\pk(x)$
              \smallskip
              \item $\pcif \ctxt = \ctxt^* :$
              \item \quad $\red{\bad \assign \true}$
              \item \quad \sout{$\Db[x] \assign \eK^*$} 
              \smallskip
              \item $\pcif \Db[x] = \bot :$ 
              \smallskip
              \item   \quad $\Db[x] \sample \bin^m$
              \smallskip
              \item $\pcreturn \Db[x]$
    				\end{nicodemus}	
    			\end{minipage}
    		}
    	\end{tabular}%
    	}
    	\end{center}
    	\caption{\label{fig:cparom-hops}
        The game hops in the proof of \cref{thm:cparom}, following
        the arrows of \cref{fig:cparom-hops-illustrated}.
        The dashed box highlights the difference between the real-world
        and the random-world indistinguishability games, the solid boxes
        highlight lines of code that were added between successive games, and
        strike-outs highlight lines of codes that were removed between games.
        Highlighted in \red{red} is the ``bad event'', represented by the flag
        $\bad$ being set to $\true$.
      }
    \end{figure}

      So, $\game 3$ will be noticeably inconsistent with the real-world game if
        $\adv$ ever calls $\RO$ on $x^*$; but what $\game 3$ \emph{will} look
        consistent with, is the case that $\adv$ was given a uniformly random
        $\eK$, chosen completely independently of the ciphertext $\ctxt^*$.
        That's the random-world game! 
        The final hop to $\game 4$ makes this manifest by removing the
        the $\bad$ flag again, as well as the re-encryption check in the
        oracle interface. Since the flag itself had no affect on the operation
        of the game, this change is once again for free, and the result is
        exactly the random-world $\indcpa$ game instantiated with our
        construction.
      \begin{align*}
        \abs{\probsample{ \game{3}(\adv) }{1 \gets \adv} 
          - \probsample{ \game{4}(\adv) }{1 \gets \adv}} = 0\, . 
      \end{align*}

      Collecting, we see that we have found the bound
      \[
        \advantage{\indcpa}{\KEM} 
            %\abs{\probsample{ \game{1}(\adv) }{1 \gets \adv} 
            %- \probsample{ \game{2}(\adv) }{1 \gets \adv}} \\
            %+ \abs{\probsample{ \game{2}(\adv) }{1 \gets \adv} 
            %- \probsample{ \game{3}(\adv) }{1 \gets \adv}}
            %+ \abs{\probsample{ \game{3}(\adv) }{1 \gets \adv}
            %- \probsample{ \game{4}(\adv) }{1 \gets \adv}} \\
            \leq 0 + \probsample{\game{3}(\adv)}{\bad} + 0 =
            \probsample{\game{3}(\adv)}{\bad}\, .
      \]

      All that remains is to bound $\probsample{\game{3}(\adv)}{\bad}$,
        and this is where our reduction $\bdv$ enters the stage.
        The reduction is playing its own $\owcpa$ game, and is given black-box
        access to the adversary $\adv$ (presumably together with high praises
        of its code-breaking abilities). It wants to leverage the abilities of
        $\adv$ to win at its own game, by simulating an environment for
        $\adv$ and extracting its outputs. However, $\bdv$ will have to
        provide an \emph{exact} simulation of the environment that $\adv$
        expects: if $\adv$ notices anything is off, it can refuse to
        participate, for example by deliberately outputting wrong values, or
        halting prematurely.\footnote{
          To see how this admittedly anthropomorphizing description arises from the
          mathematical statement, recall that the theorem quantizes
          over the distinguishing advantage of \emph{any} adversary that is
          placed in the \indcpa environment. It does not (and can not) say anything
          about how these same adversaries behave in any \emph{other}
          environments, and the ``for all'' quantifier means that we 
          might as well assume the worst.
        }

      The problem is, $\adv$ expects to be placed in the real- or random-world
        $\indcpa$ game with 50/50 probability, but $\bdv$ only knows how to
        simulate the random-world game. 
        If $\bdv$ just simulated the random-world game for $\adv$,
        then it would look to $\adv$ like $b$ were equal to $1$ with $100\%$
        probability. Hence, it would seem that an adversary that simply ignores
        its input and outputs $1$ would break $\bdv$'s strategy, since it
        would win at the simulated game $100\%$ of the time while leaving $\bdv$
        empty-handed.

      But of course, $\bdv$ will not be simulating the random-world game
        $\game 4$, but rather the equivalent $\game 3$. Notice how the $\bad$
        flag is first set to $\true$ the first time the encrypted value $x^*$
        is input to the random oracle. Notice further that, thanks to the
        re-encryption trick, this check can actually be implemented by $\bdv$,
        who knows $\pk$ and $\ctxt^*$ but not $x^*$. A strategy reveals
        itself: Use the inputs $\pk$ and $\ctxt^*$ given by the \owcpa
        environment to simulate $\game 3$ for $\adv$, except that when $\game
        3$ would set $\bad$ to $\true$, instead halt the simulation and output
        the value $x^*$ that was just entered into the random oracle.

      The pseudocode for our reduction $\bdv$ is given in
        \cref{fig:cparom-red}. Since the event that $\bdv$ halts and outputs
        the correct value corresponds exactly to the event that $\game 3$
        would set $\bad$ to $\true$, we have
      \[
        \probsample{\game{3}(\adv)}{\bad} =
          \advantage{\owcpa}{\dPKE}[(\bdv)]\, ,
      \]
      from which we may finally conclude that
      \[
        \advantage{\indcpa}{\KEM} \leq 
        \probsample{\game{3}(\adv)}{\bad}
        = \advantage{\owcpa}{\dPKE}[(\bdv)]\, . 
      \]

    %%% Reduction
    \begin{figure}[t]
    	\begin{center}
    		\fbox{
    			\small
    			\nicoresetlinenr
    			\begin{minipage}[t]{3.7cm}	
            \underline{\redHeading $\bdv(\pk, \chalctxt)$}
            \smallskip
    				\begin{nicodemus}
              \item $\eK^* \sample \Kspace$
              \smallskip
    					\item $\advbit \gets \adv^{\RO}(\pk, \eK^*, \chalctxt)$
              \smallskip
    					\item $\pcoutput \bot$
    				\end{nicodemus}
    			\end{minipage}
    
    			\begin{minipage}[t]{5cm}
    				\underline{if $\adv$ calls $\RO(x)$}
            \smallskip
    				\begin{nicodemus}
              \item $c \assign \Enc_\pk(x)$ 
              \smallskip
              \item \pcif $c = \chalctxt :$ 
              \smallskip
              \item   \quad \green{\pcabort and \pcoutput $x$}
              \smallskip
              \item $\pcif \Db[x] = \bot :$ 
              \smallskip
              \item   \quad $\Db[x] \sample \bin^m$
              \smallskip
              \item $\pcreturn \Db[x]$
    				\end{nicodemus}	
    			\end{minipage}
    		}
    	\end{center}
    	\caption{\label{fig:cparom-red}
        Pseudocode for our reduction $\bdv$;
        highlighted in green is the point at which $\bdv$ will abort the
        simulation and win the one-wayness game. 
        If the simulation runs to completion without the correct value ever having
        been called to the random oracle, $\bdv$ returns $\bot$ to indicate 
        ``I give up''.
      }
    \end{figure}

    In summary: \emph{Either} $\adv$ never calls the
      random oracle on the encrypted value $x^*$ ($\bad = \false$), 
      in which case the view of $\adv$ is
      perfectly consistent with being in either of the real- or
      random-world games (advantage $= 0$), \emph{or} $\adv$
      calls $\RO$ on $x^*$ and therefore \emph{would} be able to
      distinguish ($\bad = \true$), in which case $\bdv$ 
      aborts the simulation and outputs $x^*$, 
      winning the one-wayness game.
    \qed
  \end{proof}

  %% IN CASE OF MULTI-CHAL VERSION
    %If this were all, then we would get a completely tight bound. 
    %  There is still one avenue of attack available to $\adv$, however,
    %  namely to guess the message underlying a \emph{future} challenge,
    %  then call the random oracle on that value before $\bdv$ has a
    %  chance of recognizing it as such. When that challenge later appears,
    %  $\adv$ will notice that the key $K$ it is issued in the challenge does
    %  not match the key that it had already computed.
    %
    %In Exercise~\ref{ex:remove-loss-term}, you are tasked with coming up with
    %  a better reduction that also handles this case. 
    %  %One solution would be to let our reduction check the query database for
    %  %collisions every time it receives a challenge, which would again make for
    %  %a tight reduction. 
    %  %in order to demonstrate techniques that will be useful later on when we move to the
    %  %quantum case, 
    %  Here, we opt for a ``dumber'' route, in order to demonstrate one of our
    %  main proof techniques,
    %  namely to hop to a game that simply aborts if this ``bad event'' ever
    %  occurs, and then argue that the probability this actually makes the game abort
    %  must be small. Then we can return to our main strategy, which will now 
    %  work without issue.

\subsection{When You Have to Guess}
  What would happen if we started from a probabilistic rather than a
    deterministic PKE? Well, assuming the probability of getting the same
    ciphertext twice is small (as is usually the case), then the re-encryption
    check that the reduction $\bdv$ relied on to find the correct pre-image,
    would in all likelihood fail to identify $x^*$ even if it were input.
    In other words, the reduction no longer has the ability to \emph{recognize} $x^*$, and
    so it must resort to \emph{guessing}.

  The argument goes: If $\adv$ is able to distinguish, then it must have
    queried the random oracle on input $x^*$ at \emph{some} point. 
    Therefore, with $\bdv$ keeping track of the oracle database $\Db$
    as part of its simulation, $x^*$ must be \emph{in} $\Db$, even if $\bdv$
    is unable to tell which one it is. So it does the simplest thing it can
    think of: After $\adv$ halts and outputs a guess $\advbit$, $\bdv$ chooses
    one of the $x$-values from $\Db$ uniformly at random, and returns that one
    as its guess. Then, if $\adv$ made $q$ queries to $\RO$, and provided
    $\bad = \true$ ($\adv$ called $\RO$ on $x^*$), $\bdv$ will win the
    one-wayness game with probability $\nicefrac{1}{q}$. After a little
    manipulation of probabilities, we end up with the following bound.

  \begin{theorem}[\PKE is \owcpa $\To$ \KEM is \indcpa]
    \label{thm:cparom-guess}
    Let \PKE be a probabilistic PKE scheme, let \RO be a random oracle, 
    and let the $\KEM = \U[\PKE, \RO]$.
    Then there is a simple, fully black-box reduction $\bdv$, such that for
    every adversary $\adv$,
    %that makes at most $\qe$ calls to the challenge oracle, 
    \[
      \advantage{\indcpa}{\KEM}[\left(\adv^{\RO}\right)] 
        \leq q \cdot \advantage{\owcpa}{\PKE}[\left(\bdv^{\adv}\right)]\, .
    \]
    The adversary $\bdv$ runs in the same time as $\adv$, plus the %(relatively small) 
    overhead of lazy-sampling a random oracle.
  \end{theorem}

  %\paragraph{A Loss of Tightness.}
  In \cref{thm:cparom}, the advantage of $\adv$ was \emph{tightly} bounded by
    the advantage of $\bdv$, in the sense that the advantage of $\bdv$ had a
    factor $1$ in front of it. Now there is a factor $q$. What is the
    consequence?

  Well, say we would like the advantage of $\adv$ to be upper-bounded by
    $2^{-128}$, a common target for security. If we start from a determinstic
    PKE scheme, we can rely on \cref{thm:cparom}, and choose parameters such
    that the best-known attack against the one-wayness of \dPKE is expected
    to succeed with probability $2^{-128}$ per time step -- or equivalently,
    that the best attacks require at least $2^{128}$ time steps in expectation
    to succeed. However, with $q$ counting the total number of hash function
    evaluations, $q$ can realistically get quite large: It is for example commonly
    estimated that the hash function SHA3 has been evaluated more than
    $2^{80}$ times in the life-time operation the Bitcoin block-chain. 

  Let $2^{-\lambda}$ denote the new one-wayness advantage that our parameters
    should target, after compensating for the loss in tightness. Inserting $q
    \leq 2^{80}$ gives us:
    \[
      \advantage{\indcpa}{\KEM}[\left(\adv^{\RO}\right)] 
        \leq q \cdot 2^{-\lambda} \leq 2^{80} \cdot 2^{-\lambda} = 2^{80 - \lambda} \leq 2^{-128}
        \implies \lambda \geq 208\, .
    \]
    Therefore, if we want the distinguishing advantage to be bounded by
      $2^{-128}$, we will need parameters such that the best attacks only
      achieve a one-way advantage of $2^{-208}$. 
      Perhaps not a dramatic increase --
      increasing the sizes of keys and ciphertexts accordingly should fairly
      easily be able to accomodate this -- but experience shows
      that even a modest loss in efficiency may make our scheme look a lot
      less competitive to implementers than than it just did, whether or not
      its competitors come bundled with tightness-respecting security proofs
      of their own!

    We will repeatedly tackle the question of tightness over the course of
      this tutorial, as one of the central challenges of the QROM is to pin
      down exactly what level of tightness is achievable, and when.

  \begin{exercise}\label{ex:cparom-lossy}
    Prove \cref{thm:cparom-guess}.\\
    \emph{
      Hint: Let $\game 1$ -- $\game 4$ be as before; you only need to change how
      the $\bad$ check is performed. Write down pseudocode for a
      reduction $\bdv$ that implements the guessing strategy described above, and
      lower-bound its one-wayness advantage in terms of $\prob{\bad}$ and $q$.
    }
  \end{exercise}

\subsection{One-way to Hiding}\label{sec:C-O2H}
  \paragraph{Find-O2H.}
    \Cref{thm:cparom} may be about a specific construction (the
      \U-transform), but the take-away message is quite general: If a reduction
      has control over the random oracle, and is able to recognize the correct
      answer when it sees it, we can use the random oracle to tightly reduce
      security from relying on a stronger distinguishing-style assumption, to
      instead rely on a weaker one-wayness assumption. 
      %If the reduction further has the ability to
      %recognize the correct answer when it sees it, this can be done tightly;
      %otherwise, the reduction will come with a loss factor $q$.
      We can generalize this observation in the following ``one-way to hiding'' lemma.\footnote{
        ``Hiding'' is another name for indistinguishability-style
        security, most often seen in the context of commitment schemes, which
        are usually required to be both ``hiding'' and ``binding''.
      }
      %One for the case that the reduction can \emph{find} the value it
      %Is looking for, and one for the case that it has no choice but to \emph{guess}.

    \begin{restatable}[Find-O2H, classical]{lemma}{cOWHfindsingle}\label{lemma:co2h-find-single}
      Let $\RO$, $\GO: \bin^n \to \bin^m$ be random
      oracles such that for every $x \neq x^*$, $\RO(x) = \GO(x)$, whereas
      $\RO(x^*)$ and $\GO(x^*)$ are independently uniform. 
      Let $z$ be a value that depends arbitrarily on $\RO$, $\GO$, and $x^*$.
      Let $\indfunc^*$ be an indicator oracle that on input $x$
      returns $1$ if $x = x^*$, and otherwise returns $0$. 

      Then there is a simple, fully black-box extractor $\extractor$,
      such that for every distinguisher $\distinguisher$,
      \[
          \abs{\prob{1 \gets \distinguisher^{\RO}(z)}
              - \prob{1 \gets \distinguisher^{\GO}(z)}}
            \leq \probsample{x \gets \extractor^{\distinguisher, \indfunc^*}(z)}{x = x^*}\, .
      \]
    \end{restatable}

    It may not be clear at a glance how this lemma relates to the proof of
      \cref{thm:cparom}. After all, there there was only one random oracle in
      play, while \cref{lemma:co2h-find-single} has two. Also,
      \cref{thm:cparom} speaks of an adversary and a reduction, while
      \cref{lemma:co2h-find-single} has a distinguisher and an extractors. How
      do they connect?

    To see the connection, note how the difference between $\game 2$ and $\game 3$
      is contained to the oracle interface only. In fact, the random oracle
      in $\game 2$ and the random oracle in $\game 3$ behave identically
      relative to $\adv$'s input,
      \emph{except} on the input value $x^*$, where it returns an independently
      uniform value instead of the expected $\eK^*$. 
      If we let the random oracle in $\game 2$ be identified with $\RO$, and the
      random oracle in $\game 3$ be $\GO$, and further let $z = (\pk, \eK^*,
      \ctxt^*)$, then we can identify $\distinguisher$ with $\adv$, and things
      start to line up:
      \[
        \abs{\probsample{\game{2}(\adv)}{1 \gets \adv}
          - \probsample{\game{3}(\adv)}{1 \gets \adv}}
        = \abs{\prob{1 \gets \distinguisher^{\RO}(z)}
          - \prob{1 \gets \distinguisher^{\GO}(z)}}\, .
      \]

    As for the extractor, it is given in \cref{fig:co2h-extractor}. It should
      look familiar: aside from the use of the ``indicator oracle'' $\indfunc^*$, it is an
      almost exact adaptation of $\bdv$ from \cref{fig:cparom-red}.
      The indicator oracle is simply a formalization of the statement that ``the
      extractor is able to recognize the correct value''; we provide this
      ability in oracle form to make sure that our formalization 
      only gives the extractor the ability to recognize correct values, 
      such that it can be applied to any situation where this is the case. For
      example, in our case the $\indfunc^*$ oracle can be locally computed from
      the input $z$, via the same re-encryption trick employed by $\bdv$.
      %and not accidentally give it an enhanced ability to find $x^*$ directly.

      %%% Extractor
      \begin{figure}[t]
      	\begin{center}
      		\fbox{
      			\small
      			\nicoresetlinenr
      			\begin{minipage}[t]{3.5cm}	
              \underline{\extHeading $\extractor^{\distinguisher, \indfunc^*}(z)$}
              \smallskip
      				\begin{nicodemus}
      					\item $\advbit \gets \distinguisher^{\RO}(z)$
                \smallskip
      					\item $\pcoutput \bot$
      				\end{nicodemus}
      			\end{minipage}
      
      			\begin{minipage}[t]{5cm}
      				\underline{if $\distinguisher$ calls $\RO(x)$}
              \smallskip
      				\begin{nicodemus}
                \item \pcif $\indfunc^*(x) = 1 :$ 
                \smallskip
                \item   \quad \pcabort and \pcoutput $x$
                \smallskip
                \item $\pcif \Db[x] = \bot :$ 
                \smallskip
                \item   \quad $\Db[x] \sample \bin^m$
                \smallskip
                \item $\pcreturn \Db[x]$
      				\end{nicodemus}	
      			\end{minipage}
      		}
      	\end{center}
      	\caption{\label{fig:co2h-extractor}
          The extractor $\extractor$ of \cref{lemma:co2h-find-single}.
          Here, the extractor $\extractor$ simulates $\RO$ for the distinguisher
          $\distinguisher$; the extractor might of course just as well provide a
          simulation of $\GO$.
        }
      \end{figure}

    Taken together, we can now provide simpler proof of the crucial hop
      from $\game 2$ to $\game 3$ in \cref{thm:cparom}, by identifying $\adv$
      with $\distinguisher$, applying \cref{lemma:co2h-find-single} to get the
      bound
      \[
        \abs{\probsample{\game{2}(\adv)}{1 \gets \adv}
          - \probsample{\game{3}(\adv)}{1 \gets \adv}}
        \leq \probsample{x \gets \extractor^{\distinguisher, \indfunc^*}(z)}{x = x^*}\, ,
      \]
      then let our reduction $\bdv$ instead simulate the environment for
      $\extractor$ by sampling $\eK^*$, setting $z = (\pk, \eK^*, \ctxt^*)$, 
      use the re-encryption trick to respond to $\indfunc^*$ queries, and output
      whatever value $\extractor$ outputs once it halts. Thus,
      \[
        \probsample{x \gets \extractor^{\distinguisher, \indfunc^*}(z)}{x = x^*}
        = \probsample{x \gets \bdv^{\adv, \extractor}(\pk, \ctxt^*)}{x = x^*}
        = \advantage{\owcpa}{\KEM}[(\bdv)]\, ,
      \]
      and the theorem is proven.

  \paragraph{Guess-O2H.}
    We can likewise define a ``Guess-O2H'' lemma, in correspondence with the guessing
      strategy of \cref{thm:cparom-guess}. Here the extractor is not given
      access to an indicator oracle, making the lemma applicable to situations
      where the reduction does not have the ability to recognize the correct
      answer. Hence, Guess-O2H has a higher level of applicability than Find-O2H, at the
      cost of sacrificing tightness.

    \begin{restatable}[Guess-O2H, classical]{lemma}{cOWHguesssingle}\label{lemma:co2h-guess-single}
      Let $\RO$, $\GO: \bin^n \to \bin^m$ be random
      oracles such that for every $x \neq x^*$, $\RO(x) = \GO(x)$, whereas
      $\RO(x^*)$ and $\GO(x^*)$ are independently uniform. 
      Let $z$ be a value that depends arbitrarily on $\RO$, $\GO$, and $x^*$.

      Then there is a simple, fully black-box extractor $\extractor$,
      such that for every distinguisher $\distinguisher$ that makes at most $q$
      calls to its random oracle,
      \[
          \abs{\prob{1 \gets \distinguisher^{\RO}(z)}
              - \prob{1 \gets \distinguisher^{\GO}(z)}}
            \leq q \cdot \probsample{x \gets \extractor^{\distinguisher}(z)}{x = x^*}\, .
      \]
    \end{restatable}

    \begin{exercise}\label{ex:cparom-lossy}
      Give a simpler proof of \cref{thm:cparom-guess} by
      referencing \cref{lemma:co2h-guess-single}. \\
      \emph{
        Hint: Let $\extractor$ do the heavy lifting -- no need to even specify
        how $\extractor$ works; it is enough that the lemma states that it
        exists. Then, all that is left for $\bdv$ to do is to provide $\extractor$ with its
        expected environment.
      }
    \end{exercise}

  %\subsubsection{Multi-point O2H.}
  %  The above lemmas (\ldots)
  %  \hhnote{cont}
  %\begin{lemma}[Find-O2H, classical]\label{thm:co2h-find}
  %  let $\RO: \bin^n \to \bin^m$ be a random
  %  oracle, and let $z$ be a bitstring that depends arbitrarily on $\RO$.
  %  Let $\solset \subseteq \bin^n$ be a subset of inputs, and let $\RO
  %  \backslash \solset$ be the punctured random oracle that is equal to $\RO$
  %  on all inputs $x \not\in \solset$, but which outputs $\bot$ on all points
  %  in $\solset$. Let $\find$ be the event that the random oracle ever outputs
  %  $\bot$. Then for every distinguisher $\distinguisher$,
  %  \[
  %      \abs{\prob{1 \sample \distinguisher^{\ro}(z)}
  %          - \prob{1 \sample \distinguisher^{\ro\backslash \solset}(z)}}
  %        = \probsample{\distinguisher^{\ro \backslash \solset}(z)}{\find}\, .
  %  \]
  %\end{lemma}


  %\hhnote{
  %  forward reference to QROM proof
  %}

  %\hhnote{todo: make exercises play well with autoref}
  %\begin{exercise}\label{ex:remove-loss-term}
  %  Write an alternative reduction that gives the same bound as in
  %  \cref{thm:cparom}, except without the extra term $\nicefrac{\qe}{2^m}$
  %  in the upper bound.
  %\end{exercise}

\subsection{The Post-Quantum Setting -- What Breaks?}
  Having seen the proofs of \cref{thm:cparom} and \cref{thm:cparom-guess}, 
    %\cref{lemma:co2h-find-single} and \cref{lemma:co2h-guess-single} 
    the above ``O2H lemmas'' may seem
    almost too trivial to bother writing down. After all, the extractors in the
    lemmas do little more than lazy-sample a random oracle, then either return
    the correct value when input (if the extractor is able to
    recognize it), or output one of the queries at random (if not).
    With the analysis being a near-trivial manipulation of
    conditional probabilities, it is not clear that the modularity gained 
    outweighs the cost of switching mental models from the behaviour of adversaries
    between games, to distinguishers trying to tell oracles apart.

  However, the same can certainly not be said for \emph{quantum} random
    oracles, since now, the input to the oracle can be a quantum superposition
    of any number of values at once. How does one recognize the correct input value
    in such a case? Is the ``bad'' event we used in our analysis
    even well-defined? Can you even \emph{guess}, given that the no-cloning
    theorem bars you from copying down the input states?

  This is where the very same O2H formalism that we introduced above earns its
    keep. In fact, the O2H lemmas turns out to be one the most useful toolkits
    we have, and they show up everywhere in the QROM literature. 
    This is because they successfully shove the tricky quantum analysis
    ``under the rug'', allowing non-quantum cryptographers to design their post-quantum
    game hops without having to worry that they are making a quantum fool of
    themselves. (Or at least, worry \emph{less}.)

  We will meet our first quantum O2H
    lemmas in \cref{sec:o2h}, where we will apply them to lift
    \cref{thm:cparom} and \cref{thm:cparom-guess} to the post-quantum
    setting. But before then, it is about time we brushed up on our quantum. 

  %\todoenum[rough structure]{
  %  \item A discussion about why the previous proofs do not yield
  %    post-quantum security 
  %  \item Where the proofs break down, and why they do not admit direct
  %    lifting (not history-free)
  %}

    \clearpage
    \newcommand{\titleCcaRom}{Example 2: From CPA to CCA Security with the \U-Transform}
    \futuresect{qrot-cca-rom}{\titleCcaRom}
    \newcommand{\titleDsaRom}{Example 3: From Passive to Active Signature Security}
    \futuresect{qrot-dsa-rom}{\titleDsaRom}
    \newcommand{\titleFsRom}{Example 4: The Fiat-Shamir Transformation}
    \futuresect{qrot-fs-rom}{\titleFsRom}

  \addcontentsline{toc}{part}{Part 2. Post-Quantum Security is Not So Easy with Random Oracles}
  %\addcontentsline{toc}{part}{Part 2. Technical Difficulties in the Post-Quantum Realm}
  %\part{(Post-)Quantum Difficulties}\label{part:classical}
    \section{A Brief Quantum Primer}
  Quantum theory arose as a branch of physics, and so much of the
    standard notational conventions are hand-wavy, and often inconsistently
    implemented across the literature. For the sake of clarity and
    self-containedness, we specify all our
    conventions below, providing a brief summary of the
    building blocks that make up quantum computing theory along the way;
    for more details, see the celebrated introductory text by Nielsen and
    Chuang~\cite{NieChu10}.

  \todoenum[Provide explanations of:]{
    \item Bare-bones background on quantum states and Dirac notation; states
      vs.\ labels.
    \item $N = 2^n$; the set $[n]$; $n$-length registers are described by
      vectors in an $N$-sized Hilbert space.
    \item Greek letters: When and why I use them.
    \item Precise notational conventions regarding registers.
    \item Precise conventions regarding fonts.
    \item Notation for unitary evolutions: $\apply{U}$.
    \item Stick to pure states, and
      only mention the generalization to density matrices.
    \item Basic background on game hopping and advantages.
    \item Maybe write something about asymptotic vs.\ concrete security.
  }
  
  \paragraph{Dirac notation.}
    In general, quantum states are complex-valued vectors in a Hilbert space (a
      vector space defined over functions),
      which may or may not be finite-dimensional.\footnote{
        For example, consider the state $\ket{p(x)}$
        describing the position $x \in [0,1]$ of particle $p$; this will live
        in an infinite-dimensional Hilbert space characterized by the basis vectors
        $\{\ket{p(x)}\}_{x\in[0,1]}$.
      }
    %
    Quantum states are represented using the Dirac notation, in which a quantum state
      carrying the \emph{label} $\psi$ is written $\ket{\psi}$, its conjugate
      transpose is written $\bra{\psi} \coloneqq \ket{\psi}^\dagger \coloneqq
      (\ket{\psi}^*)^T$, and its
      inner product with itself is written $\braket{\psi|\psi} \coloneqq
      \abs{\ket{\psi}}^2$. Quantum states are unit-norm, and so we have that
      $\braket{\psi|\psi} = 1$. For two not necessarily equal quantum states
      $\ket{\psi}$ and $\ket{\varphi}$, their inner product
      $\braket{\psi|\varphi}$ is a real number between zero and one, where
      ``zero'' means the states are orthogonal, and ``one'' means the states are
      identical (modulo an unobservable global phase factor $e^{i\theta}$).

    This calculus is fine, but one may wonder, what \emph{is} a quantum
      state---in the sense, what kind of mathematical object is it? The answer
      can be found in the statement that they live in $N$-dimensional \emph{Hilbert
      space:} Hilbert spaces are vector spaces defined over \emph{functions} in
      place of vectors, where the inner product between two functions $f$ and
      $g$ on the interval $[a,b]$ is defined to be $\int_{-\infty}^{\infty}
      f(x)^* g(x) dx$. This inner product defines a norm, and a notion of
      orthogonality, from which a vector space may be constructed.
      Quantum states, then, are the functions that arise as solutions to a
      differential equation describing the dynamics of the system, i.e.
      the Schrödinger equation in the non-relativistic, non-spin case; the
      Pauli equation in the case with spin (such as electrons), and the Dirac
      equation in the relativistic case (excluding gravity). They take some
      input variable corresponding to an observable, such as the position $x$,
      and returns the \emph{amplitude} corresponding to that position; squaring
      this then gives the probability that observing the particle will yield
      $x$ as a result.

    If all of this sounds terribly complicated, then you understand why
      physicists grew so fond of Dirac's notation, which neatly swept all of
      these horrible details under the rug of conceptual clarity!

    Thankfully, things become a bit easier in the discrete case: the Hilbert
      space is now finite, and so it can be expressed in terms of a finite set of basis
      ``vectors'' (i.e., functions), and the inner
      product is now a sum rather than an integral. Further, functions now have finite domains,
      and may thus be written as $N$-length vectors instead, where each element
      corresponds to an input to the function, enabling us to
      employ our toolbox from linear algebra to study them. This is the
      approach taken in the quantum information sciences.

  \paragraph{Qubit basis states.}
    We are interested in finite-dimensional quantum states, implemented
      on a collection of $n$ qubits. A \emph{qubit} is a quantum state existing
      in the two-dimensional Hilbert space characterized by the basis states
      $\ket{0}$ and $\ket{1}$. In physical implementations, these states may for
      example be mapped to the spin of an electron (``up''/``down''), or the two lowest energy
      states of an atom, or in general any two orthogonal quantum states, where
      two quantum states are orthogonal if and only if there exists a
      measurement procedure that perfectly distinguishes them.

    It is worth emphasizing that the ``$0$'' and ``$1$'' here are \emph{labels},
      and do not say anything about the physical system in question other than that
      it contains (at least) two orthogonal quantum states, which are being
      identified with the two basis states $\ket{0}$ and $\ket{1}$.

    Since two qubit basis states $\ket{b_1}$ and $\ket{b_2}$, for $b_1, b_2 \in
      \{0,1\}$, the two states must be either identical (up to the
      aforementioned phase factor) or orthogonal, we have that
    \[
      \braket{b_1|b_2} = \delta_{b_1 b_2}\, ,
    \]
    where $\delta_{ij}$ is the Kronecker delta, equal to $1$ if $i = j$, and $0$
      otherwise.
      
  \paragraph{Combined quantum states.}
    In general, two quantum systems (i.e.\ states describing 
      distinct physical systems) may be combined via the tensor product (also
      known as the \emph{outer} product): 
      \[
        \ket{\psi} \otimes \ket{\varphi} \coloneqq \ket{\psi \varphi}\, .
      \]
  
    \noindent
    If $\ket{\psi}$ describes an $N$-dimensional quantum system, and
      $\ket{\varphi}$ describes an $M$-dimensional quantum system, then the
      combined state $\ket{\psi \varphi}$ describing the combined system will be
      $NM$-dimensional. (Note how the dimensions are \emph{multiplied} rather
      than added, as one might naïvely expect.)
  
    Applying this to qubits, given a collection of $n$ qubits $\ket{b_i}$, each initialized to
      one of the two basis states (so that $b_i \in \bin$), we may construct
      the combined state 
      \[
        \ket{b_0 b_1 \ldots b_{n-1}} \coloneqq \ket{b_0} \otimes \ket{b_1}
          \otimes \ldots \otimes \ket{b_{n-1}}\, .
      \]

    \noindent
    Since each qubit is a two-dimensional quantum state, we see that the
      resulting combined state will have dimension $N \coloneqq 2^n$. The resulting Hilbert
      space will be characterized by the computational\footnote{
        As with any vector space, Hilbert spaces permit an infinitude of basis
        choices; the ``computational'' basis is simply the name of the choice
        wherein basis states are labelled by classical bits (or, in general,
        bitstrings).
      } basis states
    \[
      \{\ket{i}\}_{\forall i \in [N]} = \{ \ket{b_0 b_1 \ldots b_{n-1}}
        \}_{\forall b_0, \ldots, b_{n-1} \in \bin}\, .
    \]

    \noindent
    This exponentiality in the dimension of quantum states as the number of
      qubits increases sits at the core of
      the quantum advantage over classical computers. As an illustrative
      example, the Hilbert space corresponding to a three-qubit quantum system
      is characterized by the eight computational basis states:
    \begin{align*}
      \{ \ket{i} \}_{\forall i \in [2^3]} 
      &= \{\ket{000},
        \ket{001},
        \ket{010},
        \ket{011},
        \ket{100},
        \ket{101},
        \ket{110},
        \ket{111}\} \\
      &= \left\{ \begin{bmatrix} 1 \\ 0 \\ 0 \\ 0 \\ 0 \\ 0 \\ 0 \\ 0 \end{bmatrix},
        \begin{bmatrix} 0 \\ 1 \\ 0 \\ 0 \\ 0 \\ 0 \\ 0 \\ 0 \end{bmatrix},
        \begin{bmatrix} 0 \\ 0 \\ 1 \\ 0 \\ 0 \\ 0 \\ 0 \\ 0 \end{bmatrix},
        \begin{bmatrix} 0 \\ 0 \\ 0 \\ 1 \\ 0 \\ 0 \\ 0 \\ 0 \end{bmatrix},
        \begin{bmatrix} 0 \\ 0 \\ 0 \\ 0 \\ 1 \\ 0 \\ 0 \\ 0 \end{bmatrix},
        \begin{bmatrix} 0 \\ 0 \\ 0 \\ 0 \\ 0 \\ 1 \\ 0 \\ 0 \end{bmatrix},
        \begin{bmatrix} 0 \\ 0 \\ 0 \\ 0 \\ 0 \\ 0 \\ 1 \\ 0 \end{bmatrix},
        \begin{bmatrix} 0 \\ 0 \\ 0 \\ 0 \\ 0 \\ 0 \\ 0 \\ 1 \end{bmatrix} \right\}\, .
    \end{align*}

    \noindent
    As in the single-qubit case, we see that we have orthonormality over the
      computational basis states:
    \[ \braket{x|y} = \delta_{xy}\, . \]

  \paragraph{Qubit registers.}
    Given a collection of $n$ qubits, we will often find it useful to divide it
      into \emph{registers} of specified sizes, for example to separate
      distincts elements in an input. We will use subscripts to refer to specific
      registers, unless the register we are referring to is clear from context.

    As an example, let's say we want the first
      $n_1$ qubits store the variable $x$---call this ``Register 1''---and that
      the remaining $n_2$ qubits---``Register 2''---should store the variable
      $y$. Then the following are all equivalent ways of denoting the combined state:

    \[
      \ket{x\concat y} = \ket{x}_1 \otimes \ket{y}_2 \coloneqq \ket{x}_1 \ket{y}_2
        \coloneqq \ket{x, y}\, ,
    \]
    \noindent
    where $\concat$ denotes string concatenation.

  \paragraph{Greek and Roman letters.}
    We will use the letters $x$, $y$, $z$, \ldots to represent classical binary
      strings. Thus, a state $\ket{x}$ will always be understood to be a
      computational basis state. Whenever we are talking specifically about single-qubit
      computational basis states, we will use $b$ to represent a single bit, 
      and write $\ket{b}$. 

    General quantum states are traditionally labelled by letters taken from the
      middle of the greek alphabet, such as $\psi$, $\varphi$, $\xi$, etc.
      We will follow this tradition, or otherwise use a descriptive label
      relating to the specific context in which the state is being used. 

    Quantum \emph{amplitudes}, meanwhile, are written using letters taken from the start of the greek
      alphabet ($\alpha$, $\beta$, \ldots). 
      
  \paragraph{General qubit states.}
    A general single-qubit state may be written in terms of the computational
      basis states as follows:

    \[
      \ket{\psi} = \alpha_0 \ket{0} + \alpha_1 \ket{1}\, ,
    \]
    where the amplitudes $\alpha_i$ are complex numbers satisfying
      $\abs{\alpha_0}^2 + \abs{\alpha_1}^2 = 1$, (where the absolute value of
      a complex number $\alpha$ is defined as $\abs{\alpha} = \sqrt{\alpha^*
      \alpha}$).

    Whenever a qubit is in a state
      described by amplitudes $\alpha_0$ and $\alpha_1$ that are both non-zero, we
      say that the qubit is in a \emph{superposition} over the values ``$0$'' and
      ``$1$''. If we then \emph{observe} the qubit, the probability that the
      observed outcome $X$ is either $0$ or $1$ is then given by $\prob{X=b} =
      \abs{\alpha_b}^2 = \abs{\braket{\psi |b}}^2$, and after observation, the superposition will have
      ``collapsed'' to the basis state corresponding to the observed value,
      $\ket{b}$.\footnote{
        The discernible effect of observations on quantum states has led to a
        century of heated philosophical debate over how quantum theory should be
        interpreted: The term ``collapse'' used here is typically associated with the
        conservative ``Copenhagen interpretation'', championed by Niels Bohr, and commonly
        characterized as taking a ``shut-up-and-calculate'' attitude towards the problem.
        Another popular interpretation is the many-worlds interpretation due to Hugh Everett,
        and favoured by 
        father of quantum computing theory David Deutsch, among others. Instead
        of the quantum state ``collapsing'' to become ``classical'',
        this interpretation would say that both terms are ill-defined, and that
        the universe should instead be understood to be nothing but one big
        superposition state over all worlds that might have arisen from
        the initial conditions of the Big Bang; and that what actually happened
        when you made your observation, was that the quantum state describing the subsystem you
        were studying and the quantum
        state describing the subsystem containing \emph{you} became correlated
        via the process known as \emph{entanglement}.
        (In case you're wondering, I place my own bet on the \emph{relational}
        interpretation due to Carlo Rovelli, which may be described as ``many-worlds
        without the multiverse''.)
      }

    Likewise, a general, $n$-qubit quantum state may be written in terms of the computational
      basis states:

    \[
      \ket{\psi} = \sum_{\forall i \in [2^n]} \alpha_i \ket{i} = \sum_{i = 0}^{N-1} \alpha_i \ket{i}
        = \alpha_0 \ket{0 \dots 00} + \alpha_1 \ket{0 \dots 01} + \ldots +
        \alpha_{N-1} \ket{1 \dots 11}\, ,
    \]
    
    \noindent
    where the \emph{amplitudes} $\alpha_i$ again satisfy
      $\sum_{\forall i \in [2^n]} \abs{\alpha_i}^2 = 1$.

    When expanded over the basis states, the inner product of two quantum
      states $\ket{\psi} = \sum_{\forall i \in [N]} \alpha_i \ket{i}$ and $\ket{\varphi} =
      \sum_{\forall j \in [N]} \beta_j \ket{j}$ may be written

    \[
      \braket{\psi|\varphi} = \sum_{\forall i, j \in [N]} \alpha^*_i \beta_j \braket{i|j} =
      \sum_{\forall i, j \in [N]} \alpha^*_i \beta_j \delta_{ij} = \sum_{\forall i \in [N]} \alpha^*_i \beta_i\, .
    \]

    \noindent
    We see that we have recovered the inner product of complex-valued vectors, which
      is just the usual inner product of real-valued vectors with an extra complex
      conjugation.\footnote{
        This is because inner products on any vector space must satisfy
        conjugate symmetry, $\langle x, y \rangle = \langle y, x \rangle^*$;
        see \url{https://en.wikipedia.org/wiki/Inner_product_space}.
      }
      Thus we may apply the usual operational interpretation of the inner
      product, in which vectors are ``perpendicular'' when their inner product is $0$,
      and ``parallel'' when it is $1$. Together with the requirement that
      quantum states have unit norm, this means that an inner product of $1$ implies
      that the two quantum states are actually identical,
      (up to the aforementioned unobservable phase factor).
      %This allows for an
      %operational interpretation of the inner product as \dots
      %$\sqrt{\abs{\braket{\psi|\varphi}}} =\sqrt{\abs{\sum_i \alpha^*_i \beta_i}}$.

    When expanding quantum states in terms of basis states in the manner above,
      we will sometimes implicitly exclude terms for which $\alpha_i = 0$,
      writing instead 

    \[
      \ket{\psi} = \sum_i \alpha_{i} \ket{x_i} \coloneqq \sum_{\forall i \in \Ispace} \alpha_{i} \ket{x_i}\, , 
    \]
    
    \noindent
    where $\Ispace = [\ell]$ indexes the set $\{\ket{x}\, |\ x \in [N]\ \land\
      \braket{\psi|x} \neq 0\}$ for some $\ell \in \NN$.
      %is the set of bitstrings such that the amplitude $\alpha_i$ of the
      %corresponding basis state $\ket{x_i}$ is non-zero.

  \paragraph{Unitary operations.}
    Since we are working over finite-dimensional Hilbert spaces, quantum
      operations, i.e.\ evolutions of quantum states, may be represented as $N
      \times N$-dimensional matrices. What properties should they have?
    
    Clearly, any operation that we perform on a quantum state (\emph{except}
      observing them!) should yield another valid quantum state
      as a result. Since quantum states must be unit-norm, this means
      that any operation we perform on them must preserve the norm. 
      This fact turns out to be sufficient\footnote{
        Here is a short proof of this fact:       } to restrict the space of allowed operations to the unitary matrices,
      i.e.\ invertible square matrices $U$ satisfying $U^{-1} = U^\dagger$.

    \begin{exercise}\label{ex:unitarity}
      Prove that requiring operations to be norm-preserving is equivalent to
      restricting the set of legal operations to the set of unitary matrices.
    \end{exercise}

    \noindent
    We will use capital roman letters to refer to quantum operations, in
      reference to their matrix representation. Whenever the qubit(s) that the
      operation is being applied to is not clear from context, we will provide
      the index of the qubit(s) in the
      subscript, e.g.\ $U_i$ for single-qubit operations, and $U_{i,j}$ for
      two-qubit operations, where $i,j \in [N]$ and $i \neq j$ (or more
      generally, the \emph{registers} on which the operations are applied).

    It is often useful to refer to operations in terms of the effect they have
      on the starting state, and so we define $\ket{\psi} \apply{U}
      \ket{\varphi}$ to mean that we perform the operation $\ket{\psi}
      \longrightarrow U \ket{\psi} = \ket{\varphi}$.

    Some elementary unitary transformations, known as quantum \emph{gates}, include
      \[
        \Xgate = \begin{bmatrix} 0 & 1 \\ 1 & 0 \end{bmatrix}\, , \quad
        \Zgate = \begin{bmatrix} 1 & 0 \\ 0 & -1 \end{bmatrix}\, , \quad
        \Hgate = \frac{1}{\sqrt{2}} \begin{bmatrix} 1 & 1 \\ 1 & -1 \end{bmatrix}\, , \quad
        \CXgate = \begin{bmatrix} 1 & 0 & 0 & 0 \\ 0 & 1 & 0 & 0 \\ 0 & 0 & 0 & 1 \\ 0 & 0 & 1 & 0 \end{bmatrix}\, .
      \]

    \begin{exercise}\label{ex:truthtable}
      Write down ``truth tables'' for each of the above gates by writing down
      their effect on the (single-qubit resp.\ two-qubit) computational basis states.
      (E.g., $\ket{0} \apply{\Xgate} \ket{1}$.)
    \end{exercise}

    \noindent
    The ``$C$'' in $\CXgate$ is short for ``controlled'', where the value of
      the first qubit \emph{controls} the operation on the second qubit. Any single-qubit
      operation $U$ defines a two-qubit controlled version of the same
      operation, which we will refer to as $CU$.

    \begin{exercise}
      Starting with the two-qubit state $\ket{\psi} = \ket{00}$, apply the operations $\Hgate_0$ and
      $\CXgate_{0, 1}$. What are the possible outcomes of observing the resulting
      state, and what are their probabilities? Compare with the result of
      applying the single-qubit gates $\Hgate_0$ and $\Hgate_1$ instead.
    \end{exercise}

    \begin{exercise}\label{ex:HadamardBasis}
      Write a new truth table for $\Xgate$, $\Zgate$, and $\Hgate$, only now in
      terms of the alternative basis states $\ket{+} \coloneqq \frac{\ket{0} + \ket{1}}{\sqrt{2}}$, 
      $\ket{-} \coloneqq \frac{\ket{0} - \ket{1}}{\sqrt{2}}$ instead. Compare
      with the truth table from \autoref{ex:truthtable}. What has changed?
    \end{exercise}

    \begin{exercise}
      Give the matrix representation of a unitary operation that acts on the
      alternative basis from \autoref{ex:HadamardBasis} in the same way that
      $\CXgate$ acts on the computational basis.
    \end{exercise}

    \clearpage
    \input{qrot-qrom}
    \clearpage
    \newcommand{\titlePseudocode}{A Quantum-Register Pseudocode}
    \section{Towards a Quantum Pseudocode}\label{sec:pseudocode}
  \subsection{Quantum Entities and Hidden Entanglement}
    In \autoref{sec:primer}, we defined notation for computational basis states
      $\ket{x}$, general quantum states $\ket{\psi}$,
      unitary transformations $\unitary$, and the
      the result of applying a unitary transformation $\unitary$ on a
      quantum state $\ket{\psi}$, written 
      $\ket{\psi} \apply{\unitary} \ket{\varphi}$; see \autoref{tab:notation}
      for an overview.

    And still this isn't enough. This is because we are going to be talking
      about communicating parties. Consider the situation that Alice has two
      quantum bits initalized to the zero state, so that her state is
      $\ket{\psi} = \ket{0}_0 \ket{0}_1$. Now she wants to send her second
      qubit register, $\ket{0}_1$, over to Bob, who is going to apply some
      single-qubit unitary to it, and then send it back. We might write this as
      $\ket{0}_1 \apply{\unitary_1} \ket{\xi}_1$, where $\ket{\xi}_1$ is some
      unknown, single-qubit state, and say that Bob returns $\ket{\xi}_1$ to
      Alice at the end of the protocol.

    But what if Alice had prepared her two qubits in the state
      $\ket{\varphi} = \frac{\ket{00} + \ket{11}}{\sqrt{2}}$ instead?
      Then, the two-qubit state is \emph{non-separable}, meaning that it can no
      longer be cleanly factored into a tensor product of two single-qubit
      states. In the world of quantum
      physics, this means that the two qubits are \emph{entangled}---a fancy
      word describing a phenomenon with many deep consequences for physics and
      philosophy alike---but mathematically speaking it is nothing but the quantum
      generalization of correlated probabilities.\footnote{
        In other words, the term ``separable'' here carries exactly the same
        meaning as the term ``product'' does when talking about classical
        probability distributions.
        More precisely, if an arbitrary quantum state is \emph{separable} along
        some dividing line, such that $\ket{\psi} = \ket{\varphi} \otimes
        \ket{\xi}$, then the probability distribution over the possible
        outcomes of observing $\ket{\psi}$ will
        form a product distribution, factorizable along the
        same dividing line. This remains true regardless of
        which basis you use to observe the state, implying that
        entanglement is an observation-independent (and therefore
        ``objective'') property of the state.
      } Alice can still send her second qubit to Bob, and Bob can still apply
      the operation $\unitary$ to it; our problem is just that with the tools
      presented so far, we have no way\footnote{
        For the readers ahead of the curve: Yes, one may reference one half of an
        entangled state using the generalized notion of a 
        \emph{mixed state}, by going to the density matrix formalization of
        quantum mechanics, which generalizes quantum states by combining them
        with classical probabilities. (Informally speaking, a density matrix is
        nothing but a probability distribution over quantum states.) The reason
        we avoid this solution is that mixed states are not
        \emph{unique}, in the sense that there are many (larger) pure
        states that might have given rise to a particular density matrix, 
        and so there would be many global states that could have led to Bob
        receiving a particular mixed state. Therefore, density matrices lack
        the expressive power to specify how Bob's operations on the qubits
        available to him may, through entanglement, affect the state of Alice's qubits.
      } of writing the operation that Bob is
      performing on $\ket{0}_1$ without referencing the entire
      state, like so:
      \[
        \frac{\ket{00}_{0,1} + \ket{11}_{0,1}}{\sqrt{2}} \longrightarrow \left( 
        \idmatrix_0 \otimes \unitary_1 \right) 
        \left( \frac{\ket{00}_{0,1} + \ket{11}_{0,1}}{\sqrt{2}} \right)\, .
      \] 

    \noindent
    What we want, is to be able to write the operations of Alice and Bob without
      having to reference qubits that they don't have access to, regardless of whether
      or not their quantum states are entangled. In other words, what we need is a
      better \emph{quantum pseudocode}.

  \subsection{Qubits and Registers}
    Let $\qubit$ be a physical qubit, in the sense of being some kind of
      particle or quantum system existing in
      the physical world. Now, guided by our intuition, we want to say that Alice,
      instead of sending a quantum \emph{state} to Bob, actually sends him the
      \emph{physical qubit} $\qubit$.\footnote{
        In real-world quantum technologies, qubits are
        typically stored in one kind of system (e.g.\ neutral atoms) and sent using another
        kind (usually photons), but due to the non-cloneability of quantum states,
        together with the universal ability to transfer any quantum state from
        one qubit to another via the $\SWAPgate$ gate, this simplified
        model is without loss of generality.
      } Anyone who receives a qubit, should also gain the
      ability to apply any transformation or measurement on it, while those who sent it 
      should lose the ability to do so, unless and until the qubit is returned to them.
      Crucially, we should be able to describe these actions \emph{locally},
      i.e.\ without reference to any qubits \emph{not} in the possession of the
      receiving party.

    So, with that goal in mind, let us introduce a little more formalism.

    \begin{definition}[Qubit]\label{def:qubit}
      A qubit $\qubit$ is a two-dimensional physical system, i.e.\ a physical
      system associated with an \emph{observable} $\hat{\qubit}$, taking values
      in $\bin$. Each qubit is additionally associated with 
      a unique \emph{register address} $i \in \NN$.\footnote{
        In the formalism of quantum Turing machines given by
        Deutsch~\cite{Deutsch85}, the addresses correspond
        to values of the \emph{address observable} $\hat{x}$
        (except that we will limit ourselves to finite quantum
        memory by letting the set of possible addresses be $[n]$
        rather than $\ZZ$).
      }     
    \end{definition}

    \noindent
    Let us refer to the qubit sitting at address $i$ as $\qubit_i$.
      If qubit $\qubit_i$ is in not entangled with any
      other qubit (i.e.\ is in a pure state), then its state may be denoted
      $\ket{\qubit_i}$ (equivalent to $\ket{\qubit_i}_i$). If
      $\ket{\qubit_i} = \ket{\psi}$ for some particular single-qubit
      state $\ket{\psi}$, then this may be denoted $\ket{\psi}_i$. Likewise,
      the state of qubit $i$ at time $t$ may be denoted $\ket{\qubit_i(t)}$.
      For the sake of illustration, let's say that at
      time $t$ the unitary $\unitary$ was applied to $\qubit_i$; this would then be written
      $\ket{\qubit_i(t)} \apply{\unitary_i} \ket{\qubit_i(t+1)}$. Likewise, if $\qubit$ was
      initialized to $\ket0$, then $\ket{\qubit(0)} = \ket0$.

    Next, let us define our ``universe'', that is, the full, undivided quantum memory 
      which will act as a stage for all our plays.

    \begin{definition}[Quantum memory]\label{def:memory}
      An $n$-qubit quantum memory is an $n$-element list of qubits $\qcomp =
      (\qubit_1, \cdots, \qubit_n)$ comprising a quantum
      system that is \emph{closed} in the following sense: 
      There are no operations that may be performed on the qubits $\qubit_i$
      such that $\qcomp$ can no longer be said to be in a pure state.
    \end{definition}

    \noindent
    In particular, $\qcomp$ can not get entangled with any
      system outside of itself---though it may be \emph{observed}, after which
      it will immediately collapse to the observed basis state.\footnote{
        Of course, from the Everettian point of view, observing a state
        \emph{is} to entangle yourself with it.
      }         

    The requirement that $\qcomp$ is always in a pure state also
      implies that state preparation must be
      \emph{deterministic} (or at least posterior to any randomness
      involved in the procedure): Flipping a (classical) coin and saying that you
      prepare a qubit in the state $\ket{1}$ on heads and $\ket{0}$ on
      tails would, when conditioning on the result of the coin flip, put $\qcomp$
      in a mixed state, and this must therefore be disallowed. Thankfully, due to
      Turing-universality, we can make this restriction without loss of generality.

    Thus, our $N$-dimensional ``global'' quantum state $\ket{\qcomp}$ will always
      be well-defined, giving us a ``smallest'' element $\qubit$ and a
      ``largest'' element $\qcomp$. That means it is time to define our in-betweens.

    \begin{definition}[Quantum register]\label{def:register}
      An $m$-qubit quantum register $\register$ is a subsystem of the $n$-qubit quantum
      memory $\qcomp$, that is, an ordered sub-list of $\qcomp$, where by
      \emph{ordered} we mean that if
      $\register = (\qubit_{i_0}, \qubit_{i_1}, \dots, \qubit_{i_{m-1}})$
      then $i_0 < i_1 < \dots < i_{m-1}$.
    \end{definition}

    \noindent
    If the qubits in $\register$ are not entangled with any other subsystems
      of $\qcomp$, then the system $\register$ is in a pure state, and we may
      write $\ket{\register}$. In a slight overload of notation, let
      $[\register]$ denote the set $\{i_j\}_{j\in [m]}$ of qubit addresses
      associated with the qubits in $\register$; then $\ket\register$ is given
      by
    \[ 
      \ket\register 
      \coloneqq \bigotimes_{\forall i \in [m]} \ket{\register[A][i]}
      \coloneqq \bigotimes_{\forall i \in [\register]} \ket{\qcomp[i]}
      = \ket{\qubit_{i_0}, \ldots, \qubit_{i_{m-1}}}\, .
    \]

    \noindent
    If $\ket{\register} = \ket{\psi}$ for some specific state $\ket{\psi}$,
      this may be denoted $\ket{\psi}_{[\register]}$,
      which, having satisfied ourselves that it is
      well-defined, we will immediately shorthand to $\ket{\psi}_{\register}$.

  \subsection{Operating on Registers}
    But of course, the starting point of this entire exercise was that we
      \emph{cannot} write
      the state of a register in terms of its own constituents in the general,
      entangled case, and so we still need an alternative way to denote 
      operating on it.
    
    \begin{definition}[Operating on registers]\label{def:apply}
      Let $\qcomp$ be an $n$-qubit quantum memory in quantum state
      $\ket{\qcomp}$, let the register $\register$ be an $m$-qubit
      subsystem of $\qcomp$, and let $\unitary$ be an
      $M \times M$ unitary matrix. 

      Next, define $\unitary_{\register}$ to be the $N \times N$ unitary matrix that
      has the effect of applying $\unitary$ on the qubits in $\register$ and the
      identity on the qubits in $\qcomp \backslash \register$. Specifically:
      \[
        \unitary_{\register} \coloneqq \SORTgate^\dagger([\register], [\qcomp
        \backslash \register])\ (\unitary \otimes \idmatrix)\
        \SORTgate([\register], [\qcomp \backslash \register])\, ,
      \]
      where $\SORTgate([\register], [\qcomp \backslash \register])$ is the $N
      \times N$ unitary matrix, composed of a series of $\SWAPgate$ gates, that 
      reorders the qubits in $\qcomp$ such that the qubits of $\register$
      precede the qubits of $\qcomp \backslash \register$;
      $\SORTgate^\dagger$ reverses the operation of $\SORTgate$; and 
      $\idmatrix$ is the $2^{n-m} \times 2^{n-m}$ identity matrix.

      Then, the application of $\unitary$ to the qubits in register $\register$ is
      captured in the method $\applyU{\unitary}{\register}[\qcomp]$, defined as
      \begin{align*}
        \applyU{\unitary}{\register}[\qcomp] &: \ket{\qcomp} \longrightarrow
          \unitary_{\register} \ket{\qcomp}\, .
      \end{align*}
    \end{definition}

    \noindent
    Whenever $\qcomp$ is clear from context, we will drop the subscript and
      simply write $\applyU{\unitary}{\register}$. 

    With this definition, it is clear that Alice \emph{may} only be operating on
      the qubits that are available to her, and yet those very same operations may
      be affecting distant parts of $\qcomp$ through entanglement, which
      is now captured in $\applyU{\unitary}{\register}[\qcomp]$ being a
      well-defined evolution of the global pure state $\ket\qcomp$.

    ubsubsection{Interacting entities.}
    With these pieces in place, all that's missing is to define what we mean
      when we say that ``Alice sends a qubit register to Bob''. We will do this
      by letting Alice and Bob each have access to the registers
      $\register[A]$ and $\register[B]$, respectively, which will themselves be
      non-overlapping subregisters of the total system $\qcomp$.

    \begin{definition}[Transferring a subregister]
      Let $\register[A]$ and $\register[B]$ be non-overlapping qubit registers. A
      subregister $\register[R]$ is said to be \emph{transferred} from
      $\register[A]$ to $\register[B]$, denoted $\register[A]
      \apply{\register[R]} \register[B]$, if the elements of $\register[R]$ are
      simultaneously removed from $\register[A]$ and added to $\register[B]$:
      \[
        \register[A] \apply{\register[R]} \register[B]
        \quad \Longleftrightarrow \quad 
        \register[A] \gets \register[A] \backslash \register[R]\, , \quad \register[B] \listapp \register[R]\, .
      \]
    \end{definition}

    \noindent
    Then, Alice sending register $\register[R]$ to Bob simply amounts to
      transferring $\register[R]$ from $\register[A]$ to $\register[B]$.

    Finally, we have reached our goal: A quantum pseudocode rich enough 
      to draw up well-defined interaction diagrams, like the one shown in
      \autoref{fig:aliceandbob}.

    \begin{figure}[t]
      \centering
      \def\arrowlength{3cm}
      \fbox{
        \pseudocode{
          \underline{\text{Alice$(\register[A])$}}\
                \<\< \underline{\text{Bob$(\register[B], \unitary)$}} \\
          (\qubit_0, \ldots, \qubit_{m-1}) \gets \register[A] \\
          \register[R] \gets (\qubit_{i_1}, \ldots \qubit_{i_\ell}) 
            \<\sendmessageright*[\arrowlength]{\register[R]} \\
            \<\< \applyU{\unitary}{\register[R]} \\
            \< \sendmessageleft*[\arrowlength]{\register[R]} \\
          x \sample \meas{}{\register[R]}
        }
      }
      \caption{
        \label{fig:aliceandbob} % *Always* put fig labels in caption, or else autoref gets confused
        Alice sends register $\register[R]$ to Bob, who applies unitary
        $\unitary$ to it and sends it back to Alice,
        who finally measures it and returns the observed value.
        The registers $\register[A]$ and $\register[B]$
        together make up the global quantum system $\qcomp$, left
        implicit in this figure.
      }
    \end{figure}

  \subsection{Measuring Registers}
    \hhnote{
      Split first part away and add it to the primer again, as
      ``Formalizing Measurements''?
    }
    \label{sec:measurements}
      Well---not quite. There is still one piece of formalism that needs to be
        dealt with before we can call it a day: Observations. Our goal is to make the $\meas{R}$
        method in \autoref{fig:aliceandbob} well defined, in much the same
        manner as we did with $\applyU{\unitary}{\register[R]}$---but first we
        will need to take a step back: For, if observations are not 
        unitary operations, what \emph{are} they?
     
    \subsubsection{Projectors.} 
      The answer may be gleaned from the fact, stated earlier, that \emph{upon
        observing a state, the quantum state collapses to the observed state}.
        This means that if you observe the state twice in a row
        (using the same measurement basis, to be defined shortly), then the
        outcome will be the same.
        The state has already collapsed; it can't collapse \emph{more}.

      If we call such an operation $\projector$, then we can write this insight as
        \[ 
          \projector \ket{\psi} = \projector (\projector \ket{\psi}) =
          \projector^2 \ket{\psi} 
          \quad \Longleftrightarrow \quad 
          \projector^2 = \projector\, .
        \]

      \noindent
      Such operators are known as \emph{projectors}, so named because they have
        the effect of \emph{projecting} a higher dimensional object onto a
        lower-dimensional hyperplane---or, in the case of global observations
        (i.e.\ measuring every qubit at once),
        all the way down to one of the basis states.

      For single-qubit states, it is easy to verify that the following
        operators satisfy this definition:

      \[
        \projector_0 \coloneqq \ket{0}\bra{0}
        = \begin{bmatrix} 1 \\ 0 \end{bmatrix} \begin{bmatrix} 1 & 0 \end{bmatrix}
        = \begin{bmatrix} 1 & 0 \\ 0 & 0 \end{bmatrix}\, , \quad
        \projector_1 \coloneqq \ket{1}\bra{1}
        = \begin{bmatrix} 0 \\ 1 \end{bmatrix} \begin{bmatrix} 0 & 1 \end{bmatrix}
        = \begin{bmatrix} 0 & 0 \\ 0 & 1 \end{bmatrix}\, . \quad
      \]

      \noindent
      From the matrix representations above, we see that effect of applying
        $\projector_0$ to a single-qubit state is to \emph{single out} the
        $\ket 0$ state from the superposition: 

      \[
        \projector_0 \ket{\psi} = \projector_0(\alpha \ket{0} + \beta \ket{1})
        = \alpha \ket{0}\, ,
      \]

      \noindent
      and applying $\projector_0$ on it again thereafter
        will have no effect, while applying $\projector_1$ will simply yield
        the zero vector. What's more, we see that 
        $\projector_0 + \projector_1 = \idmatrix[2]$. 
        Clearly, these projectors will be \emph{orthogonal} in the sense that they
        satisfy $\projector_i \projector_j = \delta_{ij} \projector_i$.

      On the other hand, these matrices are manifestly \emph{not} unitary, as
        seen from the fact that they are not even invertible. This reflects the
        irreversible loss of information induced by observation: After
        measuring a state, there is no way to recover the full pre-collapse
        state, other than to re-create it from scratch!

      Let us next define the projector $\projector_\Sspace$ onto some
        subspace $\Sspace$ of the full $N$-dimensional Hilbert space, as
        specified by a subset of the computational basis states
        $\{\ket{x_i}\}_{\forall i \in [k]}$ (where $k \leq N$). This will
        simply be the sum of the relevant basis projectors:

      \[
        \projector_\Sspace = \sum_{\forall i \in [k]} \ket{x_i} \bra{x_i} = 
        \sum_{\forall i \in [k]} \projector_{x_i}\, .
      \]

      \noindent
      The larger the sum, the ``gentler'' the measurement.

      \begin{example}\label{ex:projector}
        Let $\ket\psi = \sum_{i \in [8]} \alpha_i \ket{i}$ be a
        three-qubit state. Then the projector $\projector_{00b}$, given
        by 
        \[
          \projector_{00b} = \ket{000}\bra{000} + \ket{001}\bra{001}
        \] 
        ``selects'' the states 
        in the superposition in which the first and second qubits are zero:
        \[
          \projector_{00b} \ket\psi = \alpha_0 \ket{000} + \alpha_1 \ket{001}\, .
        \]
        We see that $\ket\psi$ has been \emph{projected} from 
        $8$-dimensional Hilbert space down to the $2$-dimensional subspace
        spanned by $\ket{000}$ and $\ket{001}$.
      \end{example}

      \noindent
      Unlike the basis-state projectors $\projector_i$, two general projectors
        $\projector_{\Sspace}$, $\projector_\Tspace$
        need \emph{not} be orthogonal, as subspaces $\Sspace$ and $\Tspace$ may
        be spanned by an overlapping set of basis vectors.

      \begin{exercise}
        Let $\projector_{0b0}$ be given by
        \[
          \projector_{0b0} = \ket{000}\bra{000} + \ket{010}\bra{010}\, ,
        \] 
        and show that $\projector_{0b0}$ and $\projector_{00b}$ are non-orthogonal,
        where $\projector_{00b}$ is as given in \autoref{ex:projector}.
      \end{exercise}

      \noindent
      We will for the most part be interested in sets of orthogonal projectors.
        From the perspective of the computational basis states $\ket x$, this just means
        that no two projectors ``accept'' the same string $x$ as part of their
        respective selections.
        %substrings of $x$ are a particular value, and that no two projectors
        %in the set checks the same bit for the same value. 
        For measurements,
        this translates to no measurement outcome being classified as ``undecided''.

      Projectors may not be unitary, but it is easy to see that they are
        \emph{Hermitian}, meaning
        that they satisfy $\projector^\dagger = \projector$. For the
        computational basis projectors, this follows trivially from the
        matrices being real and symmetric, but it also holds for general
        choices of basis states: Let $\projector_\varphi$ be the projector that
        projects a state $\ket \psi$ onto the one-dimensional subspace spanned
        by the state $\ket\varphi$. Then,
      \[
        \projector^\dagger_\varphi = (\ket{\varphi} \bra{\varphi})^\dagger 
        = \ket{\varphi} \bra{\varphi} = \projector_\varphi\, .
      \]

      \noindent
      Finally, a set of projectors $\{\projector_{\Sspace_i}\}_{\forall i \in [k]}$ is said to be
        \emph{complete} if $\sum_{\forall i \in [k]} \projector_{\Sspace_i} = \idmatrix$.
        In particular, the set of computational-basis projectors $\{
        \projector_i \}_{i \in [N]} = \{ \ket{i}\bra{i} \}_{i \in [N]}$
        is easily seen to be complete: Each element in the sum will simply add one $1$ to
        a previously-zero spot on the diagonal, and zero everywhere else; and
        the sum has $N$ elements, so the full diagonal is covered.
        In the context of measurements, completeness reflects that outcome
        probabilities should sum to one.

      \begin{exercise}
        Provide a set of projectors such that, when taken
        together with $\projector_{00b}$ from \autoref{ex:projector}, the set
        is orthogonal and complete.
      \end{exercise}

    \subsubsection{Projective measurements.}
      However, $\projector \ket{\psi}$ will not itself be a quantum
        state, since it will not have unit norm (unless $\projector =
        \idmatrix$, that is). Additionally, we have not yet
        touched on the concept of measurement \emph{outcomes}. What are the
        values that are actually observed? Something more is needed.
        %This leads us to the concept of a \emph{measurement basis}.

      Let's say that the outcome of some measurement is $m_x$, which we
        interpret as the string $x$. We can formalize this by letting
        $\measurement$ be a \emph{measurement operator} such that
        \[
          \ket\psi \apply{M} m_x \ket x\, ,
        \]
        for one of several possible strings $x$, via some as-yet-undefined
        process. We can also make it more general, by checking which of several
        subspaces $\ket\psi$ belongs to
        \[
          \ket\psi \apply{M} m_{\Sspace_i} \ket{\psi_{\Sspace_i}}\, ,
        \]
        where $\ket{\psi_{\Sspace_i}}$ is the state that results from projecting
        $\ket\psi$ onto subspace $\Sspace_i$.

        %i.e., such that the state collapses to the basis state $\ket\psi$.
      Note that this measurement operator is \emph{not} a projector, since
        repeated applications of $\measurement$ will give back increasing
        powers of $m_x$. However, provided that the measurement
        outcomes are all \emph{orthogonal} (again meaning that none of the
        possible outcomes would be classified as ``undecided''), we are
        guaranteed that there exists a complete set of orthogonal
        projectors $\{\projector_{\Sspace_i}\}_{\forall i \in [k]}$ such that 
        \[
          \measurement = \sum_{\forall i \in [k]} m_{\Sspace_i} \projector_{\Sspace_i}\, .
        \]
        In linear algebra terms, this defines the \emph{spectral decomposition} of the
          operator $\measurement$, with the eigenvalues $m_{\Sspace_i}$.

        In the case that you are just reading off the values of $n$ qubits
          directly in what we'll refer to as a ``computational-basis
          measurement'', this will be
          \[
            \measurement = \sum_{\forall i \in [N]} m_i \projector_i
            = \sum_{\forall i \in [N]} m_i \ket{i} \bra{i}\, ,
          \]
          which, when expressed in the computational basis, is just a sum of
            $N$ matrices with the single non-zero element $m_i$ in the $i$'th
            position along its diagonal.

        Finally, then, we are ready to formally define our measurement procedure.

        \begin{definition}[Projective measurement]
          Let $i$ be an $n$-length bitstring, let $\{ \projector_{\Sspace_i}
            \}_{\forall i \in [k]}$ be a complete, orthogonal set of projectors
            (projecting onto the orthogonal subspaces $\Sspace_i$ of the full Hilbert space),
            and let $\measurement \coloneqq \sum_{\forall i \in [k]} m_{\Sspace_i} \projector_{\Sspace_i}$
            be a projective measurement operator.
            %
            Then, \emph{measuring} a state $\ket{\psi}$ under $\measurement$ yields outcome
            $\Sspace_i$ with probability:
            \[
              \prob{\abs{\measurement \ket\psi} = m_{\Sspace_i}}
              = \abs{\projector_{\Sspace_i} \ket{\psi}}^2\, ,
            \]
            and, conditioned on $\Sspace_i$ being observed, the act of measuring 
            transforms the state as follows:
            \[
              \ket{\psi} \apply{\measurement} \frac{\projector_{\Sspace_i}
              \ket{\psi}}{\abs{\projector_{\Sspace_i} \ket{\psi}}}\, .
            \]
            The value $\projector_{\Sspace_i} \ket{\psi} \coloneqq \alpha_{\Sspace_i} \in \CC$
            is known as the \emph{amplitude} of observing outcome $\Sspace_i$.
        \end{definition}
        
        \noindent
        Again specializing to the case of measuring $n$ qubits in the
          computational basis and reading off the resulting string, we see that
          the probability of observing any particular string $x$ is given by
          \[
            \prob{\abs{\measurement \ket\psi} = m_x} = \abs{\projector_x \ket{\psi}}^2
            = \abs{\alpha_x \ket x}^2 = \abs{\alpha_x}^2\, ,
          \]
          with the post-measurement state 
          \[
            \ket\psi \apply{\measurement} 
            \frac{\projector_x \ket{\psi}}{\abs{\projector_x \ket{\psi}}}
            = \frac{\alpha_x \ket x}{\abs{\alpha_x \ket x}}
            = \frac{\alpha_x }{\abs{\alpha_x}} \ket x
            \longrightarrow \ket x\, ,
          \]
          as required (where the final step does away with the unobservable
          global phase $\frac{\alpha_x }{\abs{\alpha_x}} = e^{i\theta}$).
        
        You may wonder, what \emph{are} the values of the measurement outcomes, $m_x$?  Let
          us look at single-bit measurements: $m_0$ will signal outcome $\ket 0$, and
          $m_1$ will signal outcome $\ket 1$. Though tempting, we surely can't set
          $m_b = b$, since the outcome $\ket 0$ would then multiply away the
          entire state! Instead, the canonical choice\footnote{
            Historically, this choice corresponds to the possible outcomes of
            the \emph{spin} observable for fermionic particles, such as
            electrons, with the value $+1$ corresponding to
            ``spin-up'' (the spin vector points in the positive direction of
            the $z$-coordinate) and the value $-1$ meaning ``spin-down'' (the
            spin vector points in the negative $z$-direction).
          } is to let $m_0 = 1$ and $m_1 = -1$. Then, 

        \[
          \measurement = m_0 \projector_0 + m_1 \projector_1
          = \begin{bmatrix} 1 & 0 \\ 0 & -1 \end{bmatrix}\, .
        \]
        You may recognize this matrix from earlier: It's the $\Zgate$ gate! In
          fact, computational-basis measurements are often referred to as
          ``$\Zgate$ measurements'' for this reason. As you can see, this
          choice also has the useful property that the measurement operator is
          unitary, and so may be identified with an operator, with the possible
          outcomes represented by the eigenvalues and eigenstates of that operator.

        This mapping, from a unitary to a measurement, is possible to do for any
          operator that is both unitary and Hermitian, meaning that 
          $\unitary^{-1} = \unitary^\dagger = \unitary$. Such unitaries are
          easy to recognize: They are the ones
          with real values along the diagonal, and complex-conjugate symmetry
          over the off-diagonals: $\unitary{[i,j]} = \unitary{[j,i]}^*$. 

        For example, in addition to the $\Zgate$ and $\Xgate$, we also
          have a $\Ygate \coloneqq i\Xgate\Zgate$. It turns out that any
          measurement that is possible to perform on a single qubit, can be
          written as a linear combination of these three matrices, which are
          known as the \emph{Pauli matrices} after their discoverer.

        \begin{exercise}
          Show that interpreting the $\Xgate$ (a.k.a. ``NOT'') gate as a
          measurement operator leads to
          the possible outcomes $\ket{+}$ and $\ket{-}$ (where these are as
          defined in \autoref{ex:HadamardBasis}). 
          What are the values of $m_+$ and $m_-$?
        \end{exercise}

        \begin{exercise}
          Write out the matrix representation of $\Ygate$ in the computational
          basis, and find the possible measurement outcomes it gives rise to
          when interpreted as a measurement operator by identifying its
          eigenstates and eigenvalues.
        \end{exercise}

      \subsubsection{Measuring registers.}
        We now give the last piece of the puzzle, namely a pseudocode for
          perfoming measurements on registers that may or may not be entangled
          with other parts of the system (i.e.\ may or may not be \emph{pure}).

        \begin{definition}[Operating on registers]\label{def:apply}
          Let $\qcomp$ be an $n$-qubit quantum memory in quantum state
          $\ket{\qcomp}$, let the register $\register$ be an $m$-qubit
          subsystem of $\qcomp$, and let $\unitary$ be an
          $M \times M$ unitary matrix. 

          Next, define $\unitary_{\register}$ to be the $N \times N$ unitary matrix that
          has the effect of applying $\unitary$ on the qubits in $\register$ and the
          identity on the qubits in $\qcomp \backslash \register$. Specifically:
          \[
            \unitary_{\register} \coloneqq \SORTgate^\dagger([\register], [\qcomp
            \backslash \register])\ (\unitary \otimes \idmatrix)\
            \SORTgate([\register], [\qcomp \backslash \register])\, ,
          \]
          where $\SORTgate([\register], [\qcomp \backslash \register])$ is the $N
          \times N$ unitary matrix, composed of a series of $\SWAPgate$ gates, that 
          reorders the qubits in $\qcomp$ such that the qubits of $\register$
          precede the qubits of $\qcomp \backslash \register$;
          $\SORTgate^\dagger$ reverses the operation of $\SORTgate$; and 
          $\idmatrix$ is the $2^{n-m} \times 2^{n-m}$ identity matrix.

          Then, the application of $\unitary$ to the qubits in register $\register$ is
          captured in the method $\applyU{\unitary}{\register}[\qcomp]$, defined as
          \begin{align*}
            \applyU{\unitary}{\register}[\qcomp] &: \ket{\qcomp} \longrightarrow
              \unitary_{\register} \ket{\qcomp}\, .
          \end{align*}
        \end{definition}

        \begin{definition}[Measuring a register]\label{def:measure}
          Let $\qcomp$ be an $n$-qubit quantum memory in pure state
            $\ket{\qcomp}$, let the register $\register$ be an $m$-qubit
            subsystem of $\qcomp$, and let $\measurement$ be an $m$-qubit
            projective measurement that check the state for membership in
            Hilbert subspaces $\{ \Sspace_i \}$.

          Next, define $\measurement_{\register}$ to be the $n$-qubit
            projective measurement that
            has the effect of applying $\measurement$ on the qubits in $\register$ and the
            identity on the qubits in $\qcomp \backslash \register$
            (constructed in the same manner as $\unitary_{\register}$ in
            \autoref{def:apply}).

          Then, performing measurement $\measurement$ on the qubits in register
            $\register$ is captured in the (probabilistic) method
            $\meas{\measurement}{\register}[\qcomp]$ which observes
            $\ket{\qcomp}$ under $\measurement_{\register}$. Observing outcome
            $i$, designated by the measurement procedure returning value $m_i$,
            is written
          \[
            i \sample \meas{\measurement}{\register}[\qcomp]\, ,
          \]
            and the post-measurement will be 
          \[
            \ket{\qcomp} \apply{\measurement} \ket{\qcomp}_{\Sspace_i}\, ,
          \]
          where $\ket{\qcomp}_{\Sspace_i}$ is the result of projecting the
            qubits of $\register$ onto subspace $\Sspace_i$, and applying the
            identity to the rest.

          We will omit the subscript whenever $\qcomp$ is clear from context,
            and the measurement operator $\measurement$ whenever we are
            performing a computational-basis measurement; for example, measuring
            $\ell$ qubits and observing the bit string $x \in \bin^\ell$ will
            be written simply
          \[
            x \sample \meas{}{\register}\, .
          \]
        \end{definition}

        \noindent
        To showcase the utility of our notation, we can now restate the Born
          rule in a very compact form.

        \begin{definition}[The Born rule]\label{def:Born}
        If qubit register $\register$ is in pure state $\ket{\register} = \sum_i \alpha_i \ket{x_i}$, then
          \[
            \probsample{x \sample \meas{}{\register}}{x = x_i} = \abs{\alpha_i}^2\, .
          \]
        \end{definition}

    \clearpage
    %\futuresect{qrot-pseudocode}{\titlePseudocode}

  \addcontentsline{toc}{part}{Part 3. The Tools of the QRO Toolbox}
  %\part{The Tools of the Toolbox}\label{part:tools}
    \section{The O2H Toolkit}\label{sec:o2h}
  Ever since their first introduction by Unruh~\cite{EC:Unruh14}, the One-Way
    to Hiding (O2H) lemmas
    have served as indispensible tools for cryptographers seeking to prove
    security in the quantum random oracle model.
    Note the plural: Over the years, the original O2H lemma, which only
    introduced a very limited capability of reprogramming QROs, has
    been repeatedly generalized and expanded, often to different
    and seemingly-incomparable settings. Thus, there is not one
    ``multi-tool'' O2H Lemma, but several
    variants that together make up a ``toolkit''---and, as with any good
    toolkit, different tools are best fit for different jobs!

  In the following sections, we will introduce the different variants in
    increasing generality and complexity, starting with the very simplest
    variant, and showing along the way how to apply each to a security
    reduction. An overview of the O2H lemmas that will be presented is provided in
    \autoref{tab:o2h}.\footnote{
      There are other historical variants of the O2H lemma that are not covered here,
      due later variants being strict generalisations. These include (\ldots)
    }

  \begin{table}[h]
    \centering
    %
    \caption{
      \label{tab:o2h}
      An overview of the various One-way to Hiding (O2H) lemmas. 
      Here, Q $=$ Quantum-accessible oracles, M $=$ Multi-point reprogramming, A $=$ Adaptive
      reprogramming, B $=$ fully (as opposed to rewindable) Black-box access to
      the distinguisher, C $=$ Computationally (as opposed to
      statistically) hidden reprogramming, L $=$ allows Leaking information about
      reprogrammed points after reprogramming, and G $=$ Guess, meaning the
      extractor does not have access to an indicator function for the target
      value(s). $q$ bounds the total number of
      oracle calls, and $r$ bounds the number of reprogrammed points. 
      Variants in paranthesis have been superseded by later results, and are
      included mainly for historical reasons.
      Most bounds presented here are simplified and/or approximate, removing
      sums and ignoring constant factors and information-theoretic terms; for
      the exact statement, see the respective Lemma. 
      Finally, to make for easier comparison, the
      ``Set $\secpar$ to'' column gives a numerical example for how much one
      would have to increase the security parameter if one wants to bound the
      distinguishing advantage by $2^{-128}$, assuming 1) that the extractors
      advantage is bounded by $\varepsilon \leq 2^{-\secpar}$, 2) $q \leq 2^{80}$ queries
      to the random oracle (a common estimate for the total number of SHA3
      evaluations over the course of its lifetime), and 3) $r \leq 2^{10}$
      reprogrammed points (following the example of Bellare et
      al.~\cite{AC:BRTZ24}, who bounded the number of user corruptions
      similarly); the closer to the optimal value $\secpar = 128$ the better.
    }
    %
    \rowcolors{2}{}{gray!10}
    \begin{tabular}{ccccccccccc}
      \toprule
      Ref. & O2H Variant  & \rotatebox{90}{Guessing} 
                          & \rotatebox{90}{Quantum} 
                          & \rotatebox{90}{Multi-point} 
                          & \rotatebox{90}{Adaptive} 
                          & \rotatebox{90}{Computational} 
                          & \rotatebox{90}{Leaking} 
                          & \rotatebox{90}{Black-box} 
                          & Bound (simplified) & Set $\secpar$ to \\
      \midrule
      \cite{CCS:BelRog93} &
        Classical O2H (guess) 
        & \Tyes & \Tno & \Tyes & \Tyes & \Tyes & \Tyes & \Tyes &
        $\leq q \varepsilon$ & 208 \\
      \cite{CCS:BelRog93} &
        Classical O2H (search) 
        & \Tno & \Tno & \Tyes & \Tyes & \Tyes & \Tyes & \Tyes & 
        $\leq \varepsilon$ & 128 \\
      \midrule
      \cite{EC:Unruh14} &
        Original O2H 
        & \Tyes & \Tyes & \Tno & \Tno & \Tyes & \Tyes & \Tyes &
        $\lesssim q \sqrt{\varepsilon}$ & 416 \\
      \cite{C:Unruh14} &
        ``Adaptive'' O2H 
        & \Tyes & \Tyes & \Tno & \Tyes & \Tyes & \Tyes & \Tyes &
        $\lesssim q \sqrt{\varepsilon}$ & 416 \\
      \midrule
      \cite{C:AmbHamUnr19} &
        ``Semi-Classical'' O2H (guess) 
        & \Tyes & \Tyes & \Tyes & \Tno & \Tyes & \Tyes & \Tyes &
        $\lesssim q \sqrt{\varepsilon}$ & 416 \\
      \cite{C:AmbHamUnr19} &
        ``Semi-Classical'' O2H (search) 
        & \Tno & \Tyes & \Tyes & \Tno & \Tyes & \Tyes & \Tyes &
        $\lesssim \sqrt{q \varepsilon}$ & 336 \\
      \midrule
      \cite{AC:GHHM21} &
        Resampling Lemma 
        & \Tyes & \Tyes & \Tyes & \Tyes & \Tno & \Tyes & \Tyes &
        $\lesssim r \sqrt{q \varepsilon}$ & 356 \\
      %Fixed-Perm.\ RL & \Tyes & \Tyes & \Tyes & \Tno & \Tyes & \Tyes &
      %  $\advantage{\dist}{e}[\Dis{}] \lesssim \qRP \sqrt{ \qRO \cdot \advantage{\guess}{}[\Ext{}] }$ & 
      %  $\nicefrac{1}{\varepsilon} \approx \qRP^2 \qRO$ & \autoref{lemma:resampling-fp} & \cite{EPRINT:Jaeger24b} \\
      %\midrule
      \cite{PKC:PanZen24} &
        Comp.\ Ad.\ Reprog.\ Framew. 
        & \Tyes & \Tyes & \Tyes & \Tyes & \Tyes & \Tno & \Tyes &
        $\lesssim r q \sqrt{\varepsilon}$ & 436 \\
      \cite{EPRINT:Jaeger24b} &
        ``Fixed-Permutation'' O2H
        & \Tyes & \Tyes & \Tyes & \Tyes & \Tyes? & \Tyes? & \Tyes &
        $\lesssim \sqrt{q\varepsilon}$? & 336? \\
      \midrule
      \cite{TCC:BHHHP19} &
        ``Double-sided'' O2H 
        & \Tno & \Tyes & \Tno & \Tno & \Tyes & \Tyes & \Tyes &
        $\lesssim \sqrt{\varepsilon}$ & 256 \\
      \midrule
      \cite{EC:KSSSS20} &
        Meas.-Rewind-Meas. 
        & \Tno & \Tyes & \Tyes & \Tno & \Tyes & \Tyes & \Tno &
        $\lesssim q \varepsilon$ & 208 \\
      \cite{AC:GeLiaXue24} &
        Meas.-Rewind-Ext. 
        & \Tno & \Tyes & \Tyes & \Tno & \Tyes & \Tyes & \Tno &
        $\lesssim \sqrt{q} \varepsilon$ & 168 \\
      \bottomrule
    \end{tabular}
  \end{table}

  \subsection{Guess versus Find}
    Let us show that the extractor having access to both $\RO$ and $\GO$ is
      equivalent to the extractor having access to a quantum-accessible
      indicator function.
    
    \begin{lemma}[Double-sided oracle access $\Leftrightarrow$ Indicator oracle access]
      Let $\careset \subseteq \Xspace$, let $\RO, \GO: \Xspace \to \Yspace$
      be classical random oracles satisfying $\forall x \not\in \careset: \RO(x) =
      \GO(x)$, and let $\indfunc$ be an indicator oracle for $\careset$. 
      A quantum oracle algorithm with oracle access $\indfunc$ can
      faithfully simulate $\RO$ and $\GO$, and a quantum oracle algorithm with
      oracle access to $\RO$ and $\GO$ can simulate $\indfunc$, up to collisions in
      $\careset$.
    \end{lemma}

    \begin{proof}[sketch]
      First direction: The algorithm maintains two independent simulations of
      random oracles using the compressed random oracle technique; call them
      $\RO_1$ and $\RO_2$. Upon receiving a query to $\RO$, the algorithm
      responds using $\RO_1$. Upon receiving a query to $\GO$, the algorithm
      computes the output of both $\RO_1$ and $\RO_2$, storing each in separate output
      registrers, then uses $\indfunc$ to coherently implement a control for a
      CSWAP operation that swaps the first and second output register. The simulated
      $\GO$ then behaves like $\RO_1$ on inputs $x \in \careset$, and like
      $\RO_2$ on inputs $x \not\in \careset$.

      Second direction: Recall that the indicator oracle operates on an $n$-length
      input register and a single-qubit output register, which is flipped if the
      input $x$ is in $\careset$. To simulate $\indfunc$,
      the algorithm pairs the input register with an $m$-length output register
      initilazied to $\ket{0^m}$, then first calls $\RO$, followed by $\GO$. The
      resulting state will have basis states $\ket{x, \RO(x) \xor \GO(x)}$.
      Using CNOTs to implement an OR-function, the output bit is then flipped
      iff the state of the output is not equal to $\ket{0^m}$. Thus, we get an
      oracle that indicates whether $\RO(x)$ is unequal to $\GO(x)$, which is
      identical to the indicator function for $\careset$ up to collisions in
      $\careset$ (for which the simulation will return a false negative).
      Before returning the result, the state of the output register is
      uncomputed via one more call each to $\RO$ and $\GO$, removing 
      unwanted entanglement.
      \hhnote{
        To-do: Quantify both the loss due to collisions and due to any residue
        entanglement.
      }
    \end{proof}

  \subsection{Non-Adaptive State of the Art}
    In the lemmas below, $\RO$, $\GO$, $\careset$ (or $\target$), and $z$ are random
      with arbitrary joint distributions.

    \begin{lemma}[Classical O2H~\cite{CCS:BelRog93}]
      Let $\careset \subseteq \Xspace$, let $\RO, \GO: \Xspace \to \Yspace$
      be classical random oracles satisfying $\forall x \not\in \careset: \RO(x) =
      \GO(x)$, let $\indfunc$ be an indicator oracle for $\careset$, and let $z$ be arbitrary. Then 
      there are SFBB extractors $\findext, \guessext$ such that, for all distinguishers
      $\distinguisher$ making at most $\qRO$ calls to its random oracle,
      \[
        \abs{
          \prob{1 \sample \distinguisher^{\RO}(z)} - \prob{1 \sample \distinguisher^{\GO}(z)}
        } \leq \probsample{x \sample \findext^{\indfunc}(z)}{x \in \careset}\, ,
      \]
      and
      \[
        \abs{
          \prob{1 \sample \distinguisher^{\RO}(z)} - \prob{1 \sample \distinguisher^{\GO}(z)}
        } \leq \qRO \cdot \probsample{x \sample \guessext(z)}{x \in \careset}\, .
      \]
      The extractors run in time that of $\distinguisher$ plus the
      overhead of simulating at most $\qRO$ oracle calls, and
      $\findext$ additionally makes at most $\qRO$ calls to its
      indicator oracle.
    \end{lemma}

    \begin{lemma}[Original O2H Lemma~\cite{EC:Unruh14}]
      Let $\RO, \GO: \Xspace \to \Yspace$
      be random oracles satisfying $\forall x \neq \target: \RO(x) =
      \GO(x)$ (and independent otherwise), and let $z$ be a value that depends arbitrarily on $\RO, \GO$,
      and $\target$. Then there is an SFBB extractor $\extractor$ such that, for all distinguishers
      $\distinguisher$ making at most $\qRO$ calls to its random oracle,
      \[
        \abs{
          \prob{1 \sample \distinguisher^{\RO}(z)} - \prob{1 \sample \distinguisher^{\GO}(z)}
        } \leq 2\qRO \sqrt{\prob{\target \sample \extractor(z)}}\, ,
      \]
      where the extractor runs in time that of $\distinguisher$, plus the
      overhead of simulating at most $\qRO$ oracle responses.
    \end{lemma}

    \begin{lemma}[Multi-point O2H~\cite{C:AmbHamUnr19}]
      Let $\careset \subseteq \Xspace$, let $\RO, \GO: \Xspace \to \Yspace$
      be quantum random oracles satisfying $\forall x \not\in \careset: \RO(x) =
      \GO(x)$, let $\indfunc$ be a quantum-accessible indicator oracle for $\careset$, and let $z$ be arbitrary. Then 
      there are SFBB extractors $\findext, \guessext$ such that, for all distinguishers
      $\distinguisher$ making at most $\qRO$ calls to its random oracle,
      $\forall e \in \{1, \half\}$,
      \[
        \abs{
          \prob{1 \sample \distinguisher^{\RO}(z)}^e - \prob{1 \sample \distinguisher^{\GO}(z)}^e
        } \leq 2 \sqrt{
          (\qRO + 1) \cdot \probsample{x \sample \findext^{\indfunc}(z)}{x \in \careset}
        }\, ,
      \]
      and
      \[
        \abs{
          \prob{1 \sample \distinguisher^{\RO}(z)}^e - \prob{1 \sample \distinguisher^{\GO}(z)}^e
        } \leq 2 \qRO \sqrt{
          \probsample{x \sample \guessext(z)}{x \in \careset}
        }\, .
      \]
      The extractors run in time that of $\distinguisher$ plus the
      overhead of simulating at most $\qRO$ oracle responses, and
      $\findext$ additionally makes at most $\qRO$ calls to its
      indicator oracle.
    \end{lemma}

    \begin{lemma}[Single-point O2H, find~\cite{TCC:BHHHP19}]
      Let $\target \in \Xspace$, let $\RO, \GO: \Xspace \to \Yspace$
      be quantum random oracles satisfying $\forall x \neq \target: \RO(x) =
      \GO(x)$, let $\indfunc$ be a quantum-accessible indicator oracle for $\target$, and let $z$ be arbitrary. Then 
      there is an SFBB extractor $\extractor$ such that, for all distinguishers
      $\distinguisher$ making at most $\qRO$ calls to its random oracle, 
      $\forall e \in \{1, \half\}$,
      \[
        \abs{
          \prob{1 \sample \distinguisher^{\RO}(z)}^e - \prob{1 \sample \distinguisher^{\GO}(z)}^e
        } \leq 2 \sqrt{
          \prob{\target \sample \extractor^{\indfunc}(z)}
        }\, ,
      \]
      where $\extractor$ runs in time that of $\distinguisher$ plus the
      overhead of simulating at most $2\qRO$ oracle responses, and makes at most $\qRO$ calls to its
      indicator oracle.
    \end{lemma}

    \begin{lemma}[Multi-point O2H with rewinding, find~\cite{AC:GeLiaXue24}]
      Let $\careset \subseteq \Xspace$, let $\RO, \GO: \Xspace \to \Yspace$
      be quantum random oracles satisfying $\forall x \not\in \careset: \RO(x) =
      \GO(x)$, let $\indfunc$ be a quantum-accessible indicator oracle for
      $\careset$, and let $z$ be arbitrary. Then 
      there is an RBB extractor $\extractor$ such that, for all distinguishers
      $\distinguisher$ making at most $\qRO$ calls to its random oracle, 
      \[
        \abs{
          \prob{1 \sample \distinguisher^{\RO}(z)} - \prob{1 \sample \distinguisher^{\GO}(z)}
        } \leq 4 \sqrt{\qRO} \cdot \probsample{x \sample \extractor^{\indfunc}(z)}{x \in \careset}\, ,
      \]
      where $\extractor$ runs in time thrice that of $\distinguisher$ (rewinding
      twice), plus the overhead of simulating at most $6\qRO$ oracle responses,
      and makes at most $3\qRO$ calls to its indicator oracle.
    \end{lemma}

    \hhnote{
      MRM required access to both $\RO$ and $\GO$; exchanged for indicator
      function, making MRE strictly more general. Also double-check that it is
      rewinding twice.
    }

    \hhnote{Attempt at putting them all in one lemma}
    \begin{lemma}[Non-adaptive O2H, state of the art]
      Let $\RO, \GO: \Xspace \to \Yspace$
      be quantum random oracles, let $\careset \subseteq \Xspace$ be such that
      $\forall x \not\in \careset: \RO(x) = \GO(x)$, let $\indfunc$ be a
      quantum-accessible indicator oracle for $\careset$, and let $z$ be a value
      that depends arbitrarily on $\RO, \GO$ and $\careset$.
      Let $\distinguisher$ be a distinguisher that makes at most $\qRO$ calls to
      its random oracle. Then,
      \begin{enumerate}
        \item \textbf{(Classical guess)} There is an SFBB extractor $\extractor$ such that, for all $\distinguisher$ that only
          make classical calls to the random oracle, 
          \[
            \abs{
              \prob{1 \sample \distinguisher^{\RO}(z)} - \prob{1 \sample \distinguisher^{\GO}(z)}
            } \leq \qRO \cdot \probsample{x \sample \extractor(z)}{x \in \careset}\, ,
          \]
          where $\extractor$ runs in time that of $\distinguisher$ plus the
          overhead of simulating at most $\qRO$ oracle responses.
        \item \textbf{(Classical search)} There is an SFBB extractor $\extractor$ such that, for all $\distinguisher$ that only
          make classical calls to the random oracle, 
          \[
            \abs{
              \prob{1 \sample \distinguisher^{\RO}(z)} - \prob{1 \sample \distinguisher^{\GO}(z)}
            } \leq \probsample{x \sample \extractor^{\indfunc}(z)}{x \in \careset}\, ,
          \]
          where $\extractor$ runs in time that of $\distinguisher$ plus the
          overhead of simulating at most $\qRO$ oracle responses,
          and makes at most $\qRO$ calls to its indicator oracle.
        \item \textbf{(Multi-point guess)} There is an SFBB extractor $\extractor$ such
          that, for all $\distinguisher$ and $e \in \{1, \half\}$,
          \[
            \abs{
              \prob{1 \sample \distinguisher^{\RO}(z)}^e - \prob{1 \sample \distinguisher^{\GO}(z)}^e
            } \leq 2 \qRO \sqrt{
              \probsample{x \sample \extractor(z)}{x \in \careset}
            }\, ,
          \]
          where $\extractor$ runs in time that of $\distinguisher$ plus the
          overhead of simulating at most $\qRO$ oracle responses.
        \item \textbf{(Multi-point search)} There is an SFBB extractor $\extractor$ such
          that, for all $\distinguisher$ and $e \in \{1, \half\}$,
          \[
            \abs{
              \prob{1 \sample \distinguisher^{\RO}(z)}^e - \prob{1 \sample \distinguisher^{\GO}(z)}^e
            } \leq 2 \sqrt{
              (\qRO + 1) \cdot \probsample{x \sample \extractor^{\indfunc}(z)}{x \in \careset}
            }\, ,
          \]
          where $\extractor$ runs in time that of $\distinguisher$ plus the
          overhead of simulating at most $\qRO$ oracle responses,
          and makes at most $\qRO$ calls to its indicator oracle.
        \item \textbf{(Single-point search)} For $\careset = \{\target\}$, there is an SFBB extractor $\extractor$ such
          that, for all $\distinguisher$ and $e \in \{1, \half\}$,
          \[
            \abs{
              \prob{1 \sample \distinguisher^{\RO}(z)}^e - \prob{1 \sample \distinguisher^{\GO}(z)}^e
            } \leq 2 \sqrt{
              \prob{\target \sample \extractor^{\indfunc}(z)}
            }\, ,
          \]
          where $\extractor$ runs in time that of $\distinguisher$ plus the
          overhead of simulating at most $2\qRO$ oracle responses,
          and makes at most $\qRO$ calls to its indicator oracle.
        \item \textbf{(Multi-point search with rewinding)} There is an RBB
          extractor $\extractor$ such that, for all $\distinguisher$,
          \[
            \abs{
              \prob{1 \sample \distinguisher^{\RO}(z)} - \prob{1 \sample \distinguisher^{\GO}(z)}
            } \leq 4 \sqrt{\qRO} \cdot \probsample{x \sample \extractor^{\indfunc}(z)}{x \in \careset}\, ,
          \]
          where $\extractor$ runs in time at most thrice that of $\distinguisher$ (rewinding
          twice), plus the overhead of simulating at most $6\qRO$ oracle
          responses, and makes at most $3\qRO$ calls to its indicator oracle.
      \end{enumerate}
    \end{lemma}

  \subsection{Making the Lemmas Adaptive}

  % Subsection: Classical O2H
    %\section{Classical O2H and Bellare--Rogaway Encryption}
  \hhnote{
    In which I define our motivating scheme, define and prove classical
    variants of the O2H lemma, and show the classical (multi-chal.) security of BRE.
  }



  %% Extract the following into separate files and sections
      % as you get to them
  %\subsection{Single-Challenge Implication from the O2H Lemma}
  %  \hhnote{
  %    In which I introduce the original version of O2H, and use it to prove
  %    our first result: Single-challenge $\owcpa$ KEM yields $\indcpa$ PKE.
  %  }

  %\subsection{Multi-Challenge Implication from Semi-Classical O2H}
  %  \hhnote{
  %    In which I introduce the Semi-Classical O2H, and use it to prove
  %    that if we start from $\owpca$, we get something tighter than applying 
  %    the generic hybrid reduction to the previous theorem.
  %  }

  %\subsection{Tightening the Bound with the Resampling Lemma}
  %  \hhnote{
  %    In which I introduce the Resampling Lemma, and apply it to the first
  %    game hop to make it ever so slightly tighter.
  %  }

  %\subsection{Tighter Still from qPCA}
  %    \hhnote{
  %      In which I introduce qPCA, and show that if we have this, then we can
  %      exploit the modularity of the Semi-Classical O2H theorems to get a
  %      tighter result.
  %    }

  %  \paragraph{Moving between stronger notions.}
  %  \hhnote{
  %    In which I note that (with proof sketches) that what we have so far
  %    already suffices to move between CCA-notions, and even between notions
  %    with user corruptions, as long as opening challenges is disallowed,
  %    (since we can still assume that all the reprogrammed points were defined
  %    at the start of the game).
  %  }

  %\subsection{Adding Challenge Opening with Adaptive Reprogramming O2H}
  %  \hhnote{
  %    In which I introduce the notion of multi-bit, multi-challenge with sender
  %    openings, and argue that now (finally), the above techniques are
  %    insufficient. Introduce the Pan--Zeng adaptive reprogramming framework,
  %    and \emph{either} use it directly to show security, and then show how
  %    Joseph re-derived it from permutations, \emph{or} show how Joseph is able
  %    to get adaptivity from Semi-Classical O2H directly, and use this to prove
  %    the result (whichever one gives a simpler and more intuitive proof).
  %  }

    \clearpage
    \section{Example 1: One-wayness to Hiding Revisited}
 % OW -> IND example
    \clearpage
    \newcommand{\titleCompressed}{Interlude: Lazy-sampling QROs with the Compressed-Oracle Technique}
    \futuresect{qrot-compressed}{\titleCompressed}
    \newcommand{\titleExtract}{Tool 2: Learning From QRO Queries via Quantum Online-Extraction}
    \futuresect{qrot-extract}{\titleExtract}
    \newcommand{\titleCca}{Example 2: CPA to CCA Security Revisited}
    \futuresect{qrot-cca}{\titleCca}
    \newcommand{\titleAdaptive}{Tool 3: Adaptive Reprogramming with Fixed Permutations}
    \futuresect{qrot-adaptive}{\titleAdaptive}
    \newcommand{\titleDsa}{Example 3: Passive to Active Signature Security Revisited}
    \futuresect{qrot-dsa}{\titleDsa}
    \newcommand{\titleRewind}{Tool 4: Rewinding Quantum Adversaries with Measure-and-Reprogram}
    \futuresect{qrot-rewind}{\titleRewind}
    \newcommand{\titleFs}{Example 4: The Fiat-Shamir Transformation Revisited}
    \futuresect{qrot-fs}{\titleFs}

  \addcontentsline{toc}{part}{Part 4. Under the Rug}
  %\part{Under The Rug}\label{part:proofs}
    \newcommand{\titleQit}{A Brief Quantum Information Theory Primer}
    \futuresect{qrot-qit}{\titleQit}
    \newcommand{\titleProofs}{How to Prove Quantum Lemmas}
    \futuresect{qrot-proofs}{\titleProofs}

  \newpage
  \addtocontents{toc}{\protect\vspace{1.5em}}
  \addcontentsline{toc}{section}{References}
	\bibliographystyle{alpha}
  	%\bibliographystyle{splncs04}
  	%\bibliographystyle{plain}
  	\bibliography{./cryptobib/abbrev3,./cryptobib/crypto,./qrot.bib}

  \addcontentsline{toc}{part}{Appendix}
  \appendix

  \clearpage

  % Appendices input below
  \section{The Toolbox}
  Here we collect the various tools introduced over the course of the tutorial.

  \begin{table}[h!]
    \centering
    \caption{
      The QROM toolbox TEST. Security Parameter column shows what you need to
      reach $128$-bit security for $\epsilon \leq 2^{-\secpar}$ and $q \leq 2^{80}$.
    }
    \label{tab:toolbox}
    \small
    \begin{tabular}{@{}lccc@{}}
      \toprule
      Tool                  & Tightness   & Security Parameter & Statement \\
      \midrule
      \addlinespace
      \multicolumn{4}{@{}l}{\textit{Classical}} \\
      Equivalent-until-bad  & \textit{tight}           & $\secpar = 128$ & \cref{lemma:equbad} (\cpageref{sec:equbad}) \\
      %\midrule
      \addlinespace
      \multicolumn{4}{@{}l}{\textit{Programming}} \\
      O2H (Find)            & $2\sqrt{(q+1)\cdot\epsilon}$ & $\secpar = 338$  & \cref{sec:o2h-find}  \\
      O2H (Guess)           & $2q\sqrt{\epsilon}$ & $\secpar = 418$           & \cref{sec:o2h-guess} \\
      \addlinespace
      \multicolumn{4}{@{}l}{\textit{Simulation}} \\
      Compressed oracle     & \textit{tight}  & $\secpar = 128$                  & \cref{sec:cro}       \\
      \addlinespace
      \multicolumn{4}{@{}l}{\textit{Extraction}} \\
      Zhandry's extraction  & $\epsilon + \frac{1}{\sqrt{2^m}}$?  &  ? & \cref{sec:qoe} \\
      Online extraction     & (scheme-dependent)  &  ?          & \cref{sec:qoe}       \\
      \bottomrule
    \end{tabular}
  \end{table}

  \subsection{Programming Tools}
    \equbad* % Equivalent-until-bad restatement
    \cOWHfindsingle* % Classical Find-O2H, single-point
    \cOWHguesssingle* % Classical Guess-O2H, single-point

  \subsection{Simulation Tools}

  \subsection{Extraction Tools}

  \clearpage
  \iffutureversion
    \section{Quantum Pseudocode: Additional Definitions}\label{sec:pseudocode-defs}
  Here, we show how operations on qubit registers are defined for
    our quantum pseudocode, as introduced in \cref{sec:pseudocode}.

  \subsection{Operating on Registers}
    First, we formally define for what it means for a unitary to be applied to
    a qubit register.
    The main idea is that, thanks to the requirement that the global state
    $\ket{\qcomp}$ is in a well-defined pure state, we can define an operation
    on sub-registers $\register$ of $\qcomp$ in terms of a larger operation on
    the whole state.

    \begin{definition}[Operating on registers]\label{def:apply}
      Let $\qcomp$ be an $n$-qubit quantum memory in quantum state
      $\ket{\qcomp}$, let the register $\register$ be an $m$-qubit
      subsystem of $\qcomp$, and let $\unitary$ be an
      $M \times M$ unitary matrix. 

      Next, define $\unitary_{\register}$ to be the $N \times N$ unitary matrix that
      has the effect of applying $\unitary$ on the qubits in $\register$ and the
      identity on the qubits in $\qcomp \backslash \register$.       
      Then, the application of $\unitary$ to the qubits in register $\register$ is
      captured in the method $\applyU{\unitary}{\register}[\qcomp]$, defined as
      \begin{align*}
        \applyU{\unitary}{\register}[\qcomp] &\coloneqq \ket{\qcomp(t)}
        \apply{\unitary_{\register}} \ket{\qcomp(t+1)}\, .
      \end{align*}
    \end{definition}

    \noindent
    Note how local operations affecting distant qubits through entanglement
      is captured by $\applyU{\unitary}{\register}[\qcomp]$ 
      being an evolution of the global (pure) state $\ket\qcomp$.
    
    To see that $\unitary_{\register}$ is well-defined, note that it may be
      implemented as follows:
      \[
        \unitary_{\register} \coloneqq \SORTgate^\dagger([\register], [\qcomp
        \backslash \register])\ (\unitary \otimes \idmatrix)\
        \SORTgate([\register], [\qcomp \backslash \register])\, ,
      \]
      where $\SORTgate([\register], [\qcomp \backslash \register])$ is the $N
      \times N$ unitary matrix, composed of a series of $\SWAPgate$ gates, that 
      reorders the qubits in $\qcomp$ such that the qubits of $\register$
      precede the qubits of $\qcomp \backslash \register$;
      $\SORTgate^\dagger$ reverses the operation of $\SORTgate$; and 
      $\idmatrix$ is the $2^{n-m} \times 2^{n-m}$ identity matrix.

  \subsection{Sending Registers}\label{subsec:sending}
    In order to say ``Alice sends a qubit register to Bob'', we 
      give Alice and Bob to the disjoint subregisters
      $\register[A]$ and $\register[B]$, respectively, 
      of the total system $\qcomp$. 
      Then, Alice sending register $\register[R]$ to Bob simply amounts to
      transferring $\register[R]$ from $\register[A]$ to $\register[B]$.

    \begin{definition}[Transferring a subregister]
      Let $\register[A]$ and $\register[B]$ be non-overlapping qubit registers. A
      subregister $\register[R]$ is said to be \emph{transferred} from
      $\register[A]$ to $\register[B]$, denoted $\register[A]
      \apply{\register[R]} \register[B]$, if the elements of $\register[R]$ are
      simultaneously removed from $\register[A]$ and added to $\register[B]$:
      \[
        \register[A] \apply{\register[R]} \register[B]
        \quad \Longleftrightarrow \quad 
        \register[A] \gets \register[A] \backslash \register[R]\, , \quad \register[B] \listapp \register[R]\, .
      \]
    \end{definition}

  \subsection{Measuring Registers}\label{sec:measurements}
    The final piece of the puzzle is defining the measurement operation
      $\meas{}{\register}$.

    \begin{definition}[Measuring a register]\label{def:measure}
      Let $\qcomp$ be an $n$-qubit quantum memory in pure state
        $\ket{\qcomp}$, let the register $\register$ be an $m$-qubit
        subsystem of $\qcomp$, and let $\measurement$ be an $m$-qubit
        projective measurement that check the state for membership in
        Hilbert subspaces $\{ \Sspace_i \}$.

      Next, define $\measurement_{\register}$ to be the $n$-qubit
        projective measurement that
        has the effect of applying $\measurement$ on the qubits in $\register$ and the
        identity on the qubits in $\qcomp \backslash \register$

      Then, performing measurement $\measurement$ on the qubits in register
        $\register$ is captured in the (probabilistic) method
        $\meas{\measurement}{\register}[\qcomp]$ which observes
        $\ket{\qcomp}$ under $\measurement_{\register}$. Observing outcome
        $i$, designated by the measurement procedure returning value $m_i$,
        is written
      \[
        i \sample \meas{\measurement}{\register}[\qcomp]\, ,
      \]
        and the post-measurement will be 
      \[
        \ket{\qcomp} \apply{\measurement} \ket{\qcomp}_{\Sspace_i}\, ,
      \]
      where $\ket{\qcomp}_{\Sspace_i}$ is the result of projecting the
        qubits of $\register$ onto subspace $\Sspace_i$, and applying the
        identity to the rest.
    \end{definition}

    \noindent
    Finally, observe that the projective measurement $\measurement_{\register}$ can be
      constructed similarly to $\unitary_{\register}$ above.

    \clearpage
  \fi
  \section{How To Construct a Plaintext Checking Oracle}\label{sec:t-transform}
    \todoenum[Plan]{
      \item Define the T-transform.
      \item Prove the it achieves $\owpca$ in the QROM (ideally the multi-user
        setting), non-tightly from $\owcpa$ \ldots
      \item \ldots and (more) tightly from $\indcpa$.
      \item Prove that it also achieves $\owqpca$, from $\owcpa$ and from $\indcpa$.
      \item Show that this lets us put the $q_{\RO}$ loss under the square root.
    }

  \subsection{The T-transform}\label{sec:ttransform}
    \hhnote{I which I define the T-transform}

  \subsection{From CPA to PCA in the ROM}\label{sec:ttransform}
      \hhnote{
        In which I recall the lossy result from OW-CPA, and the tight
        result from IND-CPA.
      }

  \subsection{From CPA to qPCA in the QROM}\label{sec:ttransform}
    \hhnote{
      In which I show that the lossy result from OW-CPA also gives OW-qPCA with
      very few tweaks to the proof (ignoring decryption errors for the sake of
      this tutorial), and then show (using Hands-Off O2H??) that we get get
      something tighter for IND (namely only the square root loss).
    }

  %\clearpage
  \section{Select Exercise Solutions}\label{sec:solutions}
  \subsubsection{Solution to \cref{ex:unitarity}}\label{sec:solutions:unitarity}
    Let $U$ be such that, for all
      states $\ket{\psi}$, $\norm{U\ket{\psi}} = \norm{\ket{\psi}}$. Using
      the definition of the norm, this gives us that $\abs{(U\ket{\psi})^\dagger
      U\ket{\psi}} = \abs{\bra{\psi} U^\dagger U \ket{\psi}} =
      \abs{\braket{\psi|\psi}}$. Since this
      should hold for any $\ket{\psi}$, we must have that $U^\dagger U =
      \idmatrix = U U^\dagger$; the latter equality follows from the fact that,
      for all $\ket{\varphi}$ and $\ket{\xi}$, 
      $\abs{\braket{\varphi|\xi}} = \abs{\braket{\xi|\varphi}}$.
      Thus, $U$ satisfies the definition of an invertible matrix, with the
      inverse $U^\dagger$. Further, it is easy to see that the opposite also
      holds: any unitary matrix must be norm-preserving. Therefore, the
      set of norm-preserving operations on $N$-dimensional Hilbert space
      equal to the set of $N\times N$ unitary matrices.

  \subsubsection{Solution to \cref{ex:hermitian-projectors}}\label{sec:solutions:hermitian-projectors}
    Let $\projector_\varphi$ be the projector that
      projects a state $\ket \psi$ onto the one-dimensional subspace spanned
      by the state $\ket\varphi$. Then,
      \[
        \projector^\dagger_\varphi = (\ket{\varphi} \bra{\varphi})^\dagger 
        = \ket{\varphi} \bra{\varphi} = \projector_\varphi\, .
      \]

    \noindent
    Since $\ket\varphi$ is arbitrary, and any projector can be written as a sum of
      such one-dimensional projectors, \emph{and} a sum of Hermitian matrices must itself be Hermitian: 
      \[
        \projector^\dagger = \sum_i \projector^\dagger_i = \sum_i \projector_i = \projector\, ,
      \]
      this suffices to show that any projector must be Hermitian.

  %\clearpage
  %\section{Classical O2H and Bellare--Rogaway Encryption}
  \hhnote{
    In which I define our motivating scheme, define and prove classical
    variants of the O2H lemma, and show the classical (multi-chal.) security of BRE.
  }

 % Old BRE-based ROM draft

\end{document}
